%% MoCHII code paper -- draft.
%% AASTeX v7.0.2.  Target journal: ApJ / ApJS.
%% Author: Kwang-Il Seon (KASI / UST).
%%
%% Build: pdflatex mochii ; bibtex mochii ; pdflatex mochii ; pdflatex mochii
%%
%% Numerical values, validation results, and equations are documented in the
%% associated validation calculations and physics memorandum
%% (../MoCHII_v1.00/docs/MoCHII_physics.tex).

\documentclass[twocolumn,linenumbers]{aastex702}

\usepackage{amsmath}
\usepackage{amssymb}
\usepackage{booktabs}
\usepackage{array}

%% ---- convenience macros ----
\newcommand{\code}[1]{\texttt{#1}}
\newcommand{\MoCHII}{\textsc{MoCHII}}
\newcommand{\MoCafe}{\textsc{MoCafe}}
\newcommand{\Mthree}{M$^{3}$}
\newcommand{\HII}{H\,\textsc{ii}}
\newcommand{\HI}{H\,\textsc{i}}
\newcommand{\HeI}{He\,\textsc{i}}
\newcommand{\HeII}{He\,\textsc{ii}}
\newcommand{\HeIII}{He\,\textsc{iii}}
\newcommand{\Te}{T_{\mathrm{e}}}
\newcommand{\nelec}{n_{\mathrm{e}}}
\newcommand{\nH}{n_{\mathrm{H}}}
\newcommand{\xHI}{x_{\mathrm{HI}}}
\newcommand{\QH}{Q_{\mathrm{H}}}
\newcommand{\cmc}{\mathrm{cm^{3}\,s^{-1}}}

\shorttitle{MoCHII: Monte Carlo photoionization on octree AMR}
\shortauthors{Seon}

\begin{document}

\title{MoCHII: A Three-Dimensional Monte Carlo Photoionization
Code for Dusty Nebulae on Adaptive Octree Grids}

\author[orcid=0000-0001-9561-8134,gname=Kwang-Il,sname=Seon]{Kwang-Il Seon}
\affiliation{Korea Astronomy and Space Science Institute,
776 Daedeokdae-ro, Yuseong-gu, Daejeon 34055, Republic of Korea}
\affiliation{Department of Astronomy and Space Science,
University of Science and Technology, Korea,
217 Gajeong-ro, Yuseong-gu, Daejeon 34113, Republic of Korea}
\email[show]{kiseon@kasi.re.kr}

\begin{abstract}
We present \MoCHII{} (MOnte Carlo for \HII{} regions), a
three-dimensional Monte Carlo radiative-transfer code for dusty
photoionized nebulae on adaptive octree (AMR) grids.  \MoCHII{} is built
by adding a self-consistent gas-physics layer to the validated octree
transport, peel-off, and stochastic-dust modules of the dust code
\MoCafe{}, and it computes hydrogen and helium photoionization and
thermal balance, a trace-metal ionization cascade with modern analytic
atomic data, the dust and PAH continuum, and the recombination-line and
nebular-continuum diagnostics all from one radiation field.  The band
transport uses an analytic (zero-variance) direct-field estimator on the
octree together with a dedicated Cartesian digital-differential-analyzer
(DDA) backend, an explicit Monte Carlo diffuse field, dust absorption and
forward scattering, and observer peel-off imaging.  The atomic data are
fitted offline and evaluated online: exact hydrogenic \HI/\HeII{} cross
sections, Verner et al.\ (1996) fits for \HeI{} and the outer shells of
the metals joined to the Verner \& Yakovlev (1995) subshell fits above the
first inner-shell edge, hydrogenic case-B recombination from a
level-resolved cascade and Badnell
(2023) radiative and dielectronic recombination for \HeI{} and the metals,
with the ground-level rate that separates case A from case B taken from
the Milne relation on those same cross
sections and the recombination cooling taken from the same coefficients,
Voronov collisional ionization,
and charge exchange of the metals and of helium with hydrogen, so that
adding a metal ion is a data operation
rather than a code operation.  We validate the code against analytic
attenuation (rate-integral bias $+0.08\%$), the Str\"omgren sphere
(ionized-volume radius $+0.69\%$ versus an exact 1D reference; AMR identical to
a uniform grid), a thermal H/He nebula (median $0.12\%$ in $\Te$), and
the Lexington 2000 \HII{} region benchmarks.  Its line cooling is taken
from the level populations at the local electron density, the same
$n$-level solve that produces the emergent lines, so the thermal balance
and the line list are one physics.  On the two benchmarks the
$n_{\mathrm{p}}n_{\mathrm{e}}$-weighted temperature then agrees with a
\textsc{Cloudy} c25.00 calculation of the same models to $0.4\%$ at both
benchmarks, the HII40 \HeI{} front lies within $0.8\%$ in radius
of \textsc{Cloudy}'s, and the He$^{+}$/H$^{+}$ ratio and the outer radius
fall inside the published multi-code ranges at both, once the iteration is
run to a stopping test that bounds the drift still to come in hydrogen and
helium rather than the last change, and that accepts the Monte Carlo floor
only when the ionized volume has no net drift; the cooling itself is
cross-validated channel by channel against a \textsc{Cloudy} c25.00
decomposition on a fixed $(\Te, \nelec, \mathrm{ionization})$ state, where
it agrees to $1\%$.  We further demonstrate
solution-driven ionization-front re-refinement (a re-refined octree
reproduces a native high-resolution grid to $+0.012\%$ in ionized-volume radius
with $4.5\times$ fewer cells), the complete dust and PAH coupling (grain
heating energy closure $L_{\mathrm{IR}}/L_{\mathrm{abs}} = 1.000$; the
projected 8\,$\mu$m image turning from a central peak into a
limb-brightened ionization-front annulus as PAHs are destroyed in ionized
gas), and an atomic photodissociation-region
skin with photoelectric heating.  \MoCHII{} is written for modern
three-dimensional synthetic observations of \HII{} regions and their
surroundings.
\end{abstract}

% The following project overview is retained in source for reference but is
% omitted from the manuscript; the scholarly comparison is in
% Section~\ref{sec:compare}.
\iffalse
\section{Comparison with existing MCRT photoionization codes}

Recent developments broaden the Monte Carlo photoionization landscape.
COLT has added a coupled ionization and thermal-equilibrium workflow for
nebular emission from galaxy simulations; SIROCCO provides Monte Carlo
ionization and resonance-line transfer for compact-object outflows; and the
new ionized-gas module in SKIRT couples three-dimensional MCRT to Cloudy
tables and an inline ion solver.  These are important complementary efforts,
but they do not remove the specific design advantages of MoCHII for dusty,
structured H\,\textsc{ii} regions.

MoCHII is a self-contained chain from the ionizing source to the gas state,
dust energy budget, line emissivity, and observer product.  Its
photoionization cross sections, recombination and charge-exchange fits,
Tier-1 cooling fits, Tier-2 collision-strength fits, recombination lines,
nebular continuum, and dust coupling are stored as provenance-bearing local
data.  This differs from a hybrid workflow in which gas temperature and
opacity are supplied by precomputed Cloudy tables.  A discrepancy can thus
be traced to a particular rate, fit, or solver instead of an opaque table
interpolation.

The spatial combination is also distinctive.  MoCHII accepts simulation-
derived three-dimensional density fields, uses octree AMR, and can rebuild
the tree from the computed ionization front and clump surfaces.  The recorded
refinement gate reproduces the native high-resolution Str\"omgren radius to
$0.012\%$ with $4.5$ times fewer leaves.  Its memory-minimal uniform
Cartesian backend and incremental DDA provide a complementary path when a
tree is unnecessary.

MoCHII treats gas and dust as an explicit energy ledger rather than as an
attenuation-only post-processing step.  Frequency-dependent FUV/EUV
absorption, dust scattering, live dust and PAH survival, equilibrium grain
temperature, SEDust stochastic emission, and leaf dust-band emissivities all
use the same gas state.  The integrated IR luminosity is constrained to the
absorbed dust power, and 8, 24, 70, and 160-$\mu$m maps are spatially
resolved on the AMR mesh.

The observer layer is another MoCHII strength.  HDF5/FITS emissivity files
retain the leaf state required for map reconstruction; the Python reader
deposits AMR leaves with exact area overlap; and peel-off separates direct,
diffuse, and dust-scattered contributions for multiple observers.  The direct
thin-field gate agrees with the analytic attenuated flux to $+0.017\%$.
This is stronger than a generic projection of intrinsic cell emissivities
when escape fractions, viewing angles, or dust-scattered light are scientific
observables.

Finally, MoCHII exposes both cell-max and volume-integrated convergence
measures and maintains separate gates for transport, one-dimensional
references, atomic fits, AMR, dust, IR closure, peel-off, and map
conservation.  Its stellar and diffuse launches use separate scrambled Sobol streams
by default. This randomized quasi-Monte Carlo sampling can
reduce direct-field noise while variable-length scattering remains
pseudo-random, and global photon indexing keeps the launch set independent
of MPI rank count.

The comparison remains scoped.  COLT is stronger for galaxy-scale nebular
post-processing, SIROCCO for velocity-dependent resonance-line transfer in
compact-object winds, and SKIRT for a mature general-purpose dust-MCRT
ecosystem.  MoCHII's distinctive advantage is the combination of a
self-contained auditable gas solver, solution-driven 3D AMR, closed dust and
IR energy accounting, direct/diffuse/scattered observer estimators, and
conservation-tested resolved products in one H\,\textsc{ii}-region workflow.
\fi

\keywords{\uat{H II regions}{694} --- \uat{Interstellar medium}{847}
--- \uat{Radiative transfer}{1335} --- \uat{Photoionization}{2060}
--- \uat{Interstellar dust}{836} --- \uat{Astronomy software}{1855}}

%==========================================================================
\section{Introduction}
\label{sec:intro}

Photoionized regions --- \HII{} regions, planetary nebulae, the diffuse
ionized medium, and the circumgalactic gas illuminated by stars and
active nuclei --- are among the most information-rich objects in
astrophysics.  Their emission-line and continuum spectra encode the
ionizing radiation field, the elemental abundances, the electron
temperature and density, and the dust content of the gas.  Extracting
that information requires a radiative-transfer model that couples the
ionizing radiation field to the ionization, thermal, and dust state of
the gas self-consistently.

The classic tools for this problem are one-dimensional: \textsc{Cloudy}
\citep{cloudy2017,cloudy2023} and \textsc{Mappings}
\citep{sutherland1993,mappingsv2017} solve the coupled ionization,
thermal, and (in the case of \textsc{Mappings}) shock structure of a
spherical or plane-parallel slab with very detailed microphysics: large
ion inventories, density-dependent multilevel cooling, non-equilibrium
and time-dependent options, and grain and PAH physics.  Their strength is
completeness of the local physics; their limitation is geometry.  Real
\HII{} regions are clumpy, aspherical, and shadowed, and the diffuse
recombination field they produce is intrinsically three-dimensional.

Three-dimensional Monte Carlo photoionization codes were developed to
address exactly this geometry.  \textsc{Mocassin}
\citep{ercolano2003,ercolano2005,ercolano2008} pioneered the
self-consistent Monte Carlo treatment of the direct and diffuse fields on
a Cartesian grid, with a large ion and dust inventory, and it defines the
\emph{de facto} three-dimensional benchmark suite (the Lexington/Meudon
comparison cases; \citealt{ferland1995,pequignot2001}).  The code of
\citet{wood2004} demonstrated Monte Carlo photoionization with an
efficient path-length estimator for the mean intensity.  \textsc{M}$^{3}$
(Messenger Monte Carlo \textsc{Mappings V}; \citealt{jin2022,jin2022fractal})
grafts three-dimensional Monte Carlo transport onto the extensive
\textsc{Mappings V} microphysics, combining broad atomic physics with
arbitrary Cartesian geometry.  \textsc{CMacIonize}
\citep{cmacionize2018} adds a task-based, unstructured (Voronoi or
adaptive Cartesian) grid and a modern software design.  These codes
share a common architectural constraint: their spatial grids are either
uniform Cartesian or, in the \textsc{CMacIonize} case, unstructured, and
their diffuse-field and grid-traversal designs are tightly coupled.

\MoCHII{} (MOnte Carlo for \HII{}
regions)\footnote{\url{https://github.com/seoncafe/MoCHII}} is a new
three-dimensional
Monte Carlo photoionization code built around a different design choice:
it reuses the mature octree adaptive-mesh-refinement (AMR) engine,
event-driven traversal, peel-off imaging, and stochastic dust/PAH
emission of the dust radiative-transfer code \MoCafe{} (Seon, in
preparation)\footnote{\url{https://github.com/seoncafe/MoCafe}}, and
adds a gas-physics layer on top.  The motivation is
practical.  Retrofitting an octree traversal onto an existing uniform-grid
photoionization code is difficult: the ray-segment bookkeeping is
entangled with the physics loop, and structural rewrites are error-prone
(the \textsc{Mocassin} \code{pathSegment} traversal, for instance,
carries a dual position representation with many mutation sites).  An
event-driven AMR transport engine, by contrast, already \emph{is} the
design such a retrofit would have to reach.  \MoCafe{} provides that
engine, validated for dust transport, so \MoCHII{} inherits octree AMR,
an analytic direct-field estimator, a dedicated Cartesian
digital-differential-analyzer (DDA) backend, peel-off synthetic imaging,
MPI-3 shared-memory parallelism, and HDF5/FITS I/O without modifying the
dust tree.

On top of that engine \MoCHII{} adds: (i) a continuous-energy ionizing
photon sampler and a zero-variance direct-field estimator;
(ii) hydrogen and helium photoionization and thermal balance with modern
recombination data; (iii) a trace-metal ionization cascade organized as a
data-driven registry, so that adding an ion (currently eleven elements,
C through Ca) is an offline data operation; (iv) an explicit Monte Carlo
diffuse field with a case-A/B option; (v) the complete dust and PAH
coupling in the ionizing and far-ultraviolet (FUV) bands, with the
transport cross sections and the emission taken from one grain model and
the reradiated infrared transported in turn; (vi) solution-driven
ionization-front re-refinement of the octree; and (vii) recombination-line,
collisional-line, and nebular-continuum diagnostics with leaf-resolved
emissivity output for synthetic imaging.  The atomic data are fitted
offline from CHIANTI and the Strathclyde/Badnell tables and evaluated by
analytic formulae at run time, with a provenance header on every
coefficient file.

The remainder of this paper is organized as follows.
Section~\ref{sec:compare} compares \MoCHII{} with the established 1D and
3D codes.  Section~\ref{sec:methods} describes the Monte Carlo transport,
the grids, the diffuse field, the iteration and convergence scheme, and
the quasi-random launch option.  Section~\ref{sec:atomic} documents the
atomic data and how they are fitted offline.
Section~\ref{sec:validation} presents the validation against analytic
solutions and the Lexington benchmarks.
Section~\ref{sec:refine} demonstrates the solution-driven I-front
re-refinement and a cosmological post-processing application.
Section~\ref{sec:dust} presents the dust and PAH coupling, and
Section~\ref{sec:pdr} the photodissociation-region skin.
Section~\ref{sec:summary} summarizes and outlines planned extensions.

%==========================================================================
\section{Comparison with existing codes}
\label{sec:compare}

Table~\ref{tab:compare} places \MoCHII{} next to the reference codes along
the axes most relevant to three-dimensional nebular modeling.  The
comparison is architectural rather than a ranking of the codes:
\textsc{Cloudy} and \textsc{Mappings} carry more complete local
microphysics (larger ion inventories, full multilevel cooling of every ion,
non-equilibrium and shock options), and \textsc{Mocassin} carries a
comparably broad ion and dust inventory in three dimensions.  \MoCHII{}'s
distinguishing features are the octree AMR grid with solution-driven
re-refinement, the reuse of a validated dust-transport engine, which couples
dust absorption, forward scattering, equilibrium and stochastic emission,
and PAH physics to the gas solution, and the
peel-off synthetic-imaging output.

The grid axis is the sharpest distinction.  \textsc{Mocassin} and
\textsc{M}$^{3}$ use uniform Cartesian grids; \textsc{CMacIonize} uses an
unstructured Voronoi or adaptive Cartesian grid; \MoCHII{} uses an octree
whose refinement can be driven by the ionization solution itself
(Section~\ref{sec:refine}).  On the Str\"omgren problem this lets
\MoCHII{} concentrate resolution at the ionization front and reproduce a
native high-resolution run with several times fewer cells.

The diffuse-field axis separates the on-the-spot (OTS) approximation from
explicit transport.  \MoCHII{} defaults to an explicit Monte Carlo diffuse
field with case-A rates and ground-recombination emission, matching the
\textsc{Mocassin} and \textsc{Cloudy} treatment, and the benchmark runs of
Section~\ref{sec:hii} use it, so that the transfer treatment is the same on
both sides of that comparison.  Case-B OTS (no diffuse packets, consistent
with the analytic direct-field estimator) is the alternative, in which the
helium recombination photons that can ionize hydrogen are given to it where
they are made.

The atomic-data axis reflects vintage.  \MoCHII{} uses exact hydrogenic
cross sections, the Verner et al.\ (1996) and Verner \& Yakovlev (1995)
fits combined as Verner's own \code{phfit2} combines them, a
level-resolved hydrogenic cascade for H$^{+}$ and He$^{2+}$ and Badnell (2023)
recombination for He$^{+}$ and the metals,
CHIANTI v11 collision strengths, and a ground-level recombination
coefficient obtained from those cross sections through the Milne relation
--- the same modern data, and the same construction of the case A/B split
(a level-resolved sum and its ground term),
that \textsc{Cloudy} uses --- rather than the older compilations still
carried by \textsc{Mocassin}.

The metal-treatment axis is easy to state wrongly, so we state it from the
distributed sources.  It is tempting to say that the established codes solve
one coupled nonlinear system over every element while \MoCHII{} treats the
metals as a perturbation.  That is not what the codes do.  None of
\textsc{Mocassin}, \textsc{CMacIonize}, \textsc{Torus}, the code of
\citet{wood2004}, or even \textsc{Cloudy} assembles the elements into a joint
system; all five advance one element at a time inside nested outer
iterations, with $T$ an outer root find in every case.  Four of the five use
the same nearest-neighbor ratio chain \MoCHII{} uses, adjacent-stage ratios
$f_{i+1}/f_i$ chained and normalized.  \textsc{Cloudy} does build a matrix,
but one per element, and needs it because Auger ejection jumps several stages
at once, which no adjacent-stage chain can represent.  The ratio chain is
therefore the standard method of the field rather than a \MoCHII{}
approximation, and these codes differ not in how the stage populations are
obtained but in which feedbacks are wired in.

On this axis, \MoCHII{} includes metal photoionization opacity in the
attenuation of the ionizing field
(\code{ion\_metal\_abs}), where \textsc{CMacIonize}, \textsc{Torus}, and
\citet{wood2004} accumulate metal photoionization rates while never removing
those photons from the beam, so that a metal photoionization costs the
radiation field nothing.  That is not a small choice: switching the opacity
off moves $r(\mathrm{He}^{0}\!=\!0.5)$ inward by 0.615~pc at HII40, because
the stages with thresholds between 21 and 40~eV preferentially remove the
softest helium-ionizing photons while $\sigma(\HeI)/\sigma(\HI)$ rises from
6.1 at 25~eV to 19.9 at 95~eV, so the hardened surviving field is taken up by
$\mathrm{He}^{0}$ rather than $\mathrm{H}^{0}$.  \MoCHII{} also solves the
metal contribution to $\nelec$ as a fixed point rather than a single
evaluation (\code{metal\_ne}), since the stage fractions depend on $\nelec$
themselves and one evaluation at the H/He-only electron density is the
$\nelec \to 0$ limit in which nothing recombines; \textsc{CMacIonize} and
\citet{wood2004} count no metal electrons at all.  \code{metal\_heat} adds
metal photoheating, which \textsc{Mocassin} includes and those two omit,
together with the two losses that close the same energy budget, the
kinetic energy removed by metal recombination and the ionization energy
removed by collisional ionization (Section~\ref{sec:recomb}).
Both terms are enabled in the benchmark calculations of
Section~\ref{sec:hii}.  The trace separation describes the solution method,
rather than an omission of these feedbacks.

One feedback the other four Monte Carlo codes also omit is charge exchange
being two-way in the bookkeeping:
$\mathrm{X}^{i+} + \mathrm{H}^{+} \to \mathrm{X}^{(i+1)+} + \mathrm{H}^{0}$
recombines the metal and ionizes hydrogen in the same event, so it is a source
and a sink on both sides, and \textsc{Cloudy} alone puts hydrogen's side into
hydrogen's own balance.  \MoCHII{} includes this coupling
(\code{par\%charge\_exchange}); a fixed-state validation holds all 29
reactions to $6.5\times10^{-16}$ between the two sides.  The correction is
small here for a structural reason: where the near-resonant
$\mathrm{O}^{0}/\mathrm{H}^{+}$ pair locks oxygen to hydrogen, the two terms
cancel identically in hydrogen's balance, so what survives is the departure
from that lock.  Measured,
$\langle T[N_{\mathrm{p}}N_{\mathrm{e}}]\rangle$ moves $-1.4$~K at HII40 and
$-0.1$~K at HII20, and every front by under $0.001$~pc --- 58 and 533 times
inside the tolerance these runs were converged to.
The same near-cancellation dictates how the balance is solved: it leaves the
substitution map marginally stable in partially neutral gas, so \MoCHII{}
solves hydrogen's balance for its root rather than substituting into it, with
the cascade re-evaluated at each trial.  The unknown is
$u = \ln(n_{\mathrm{HII}}/n_{\mathrm{HI}})$, from which both fractions follow
to full relative precision, so neither can leave the physical range and the
balance keeps its digits beyond the front, where $x_{\mathrm{HII}}$ falls to
$10^{-10}$ and below; one cell solve
converges to better than $10^{-12}$ where substitution left $2\times10^{-3}$, and
the residual beyond the front falls three orders of magnitude while
$\langle T[N_{\mathrm{p}}N_{\mathrm{e}}]\rangle$ moves 0.01~K.  The one
remaining absence is metal recombination continua as a diffuse source, worth
0.29\% of the explicit diffuse channel that dominates the helium budget; it
and the measurements that size it are planned in
\code{md/METAL\_COUPLING\_PLAN.md}, and the full source-level comparison is
in \code{docs/MoCHII\_code\_comparison.tex}.  Charge exchange between
helium and hydrogen, which \textsc{Cloudy} and \textsc{Mocassin} carry for
$\mathrm{He}^{+} + \mathrm{H}$, is included in both balances in the same
two-way form (Section~\ref{sec:ci}).  What the registry buys is
unchanged: adding an ion remains a data-file addition rather than a change to
a coupled solver (Section~\ref{sec:atomic}).

A related but separate axis is the cooling inventory --- which channels a
code carries at all, since that is what sets the equilibrium temperature.
As implemented in the distributed source, the Monte Carlo code of
\citet{wood2004} cools by ten five-level metal ions (N\,\textsc{i},
N\,\textsc{ii}, O\,\textsc{i}, O\,\textsc{ii}, O\,\textsc{iii},
Ne\,\textsc{iii}, S\,\textsc{ii}, S\,\textsc{iii}, C\,\textsc{ii},
C\,\textsc{iii}) plus two two-level ions ([N\,\textsc{iii}]
57\,$\mu$m and [Ne\,\textsc{ii}] 12.8\,$\mu$m), and \textsc{CMacIonize} by
the same ten five-level ions plus three two-level ions, adding
[S\,\textsc{iv}] 10.5\,$\mu$m; both then add free-free emission and
hydrogen and helium recombination cooling, in each case with the H$^{+}$
and He$^{+}$ terms only.  \MoCHII{}
carries line cooling for 28 ions --- \HI{} and the 27 metal ions of the
eleven-element registry (Appendix~\ref{app:atomic}) --- and adds three
channels that neither of those two codes carries: \HI{} collisional
excitation, collisional-ionization cooling of H and He, and the He$^{++}$
terms of the free-free and recombination sums.  It further carries the
metal recombination and collisional-ionization losses, and the
atomic-hydrogen impact on the ground-term fine structure of C\,\textsc{i},
C\,\textsc{ii} and O\,\textsc{i}, solved in one statistical equilibrium
with the electron impact (Section~\ref{sec:cooldata}).  In a $10^4$~K \HII{}
region these are small (\HI{} collisional lines are $0.7\%$ of the
HII40 cooling budget), but they are the terms that take over in the hot,
partly neutral gas of an ionization front and in the shielded gas beyond
it.

\begin{deluxetable*}{>{\raggedright\arraybackslash}p{2.4cm}>{\raggedright\arraybackslash}p{2.5cm}>{\raggedright\arraybackslash}p{2.3cm}>{\raggedright\arraybackslash}p{2.3cm}>{\raggedright\arraybackslash}p{2.3cm}>{\raggedright\arraybackslash}p{2.3cm}}
\tablecaption{Architectural comparison of \MoCHII{} with reference
photoionization codes.  ``OTS'' = on-the-spot (case B); ``explicit'' =
transported diffuse field.  This table compares design, not the breadth
of local microphysics, on which the 1D codes and \textsc{Mocassin} are
more complete.  The two metal rows are read from the distributed sources
for \MoCHII{}, \textsc{Mocassin}, \textsc{CMacIonize}, and \textsc{Cloudy};
the \textsc{M}$^{3}$ entries follow its documentation.\label{tab:compare}}
\tablewidth{0pt}
\tabletypesize{\footnotesize}
\tablehead{\colhead{} & \colhead{\MoCHII{}} & \colhead{MOCASSIN} &
\colhead{M$^{3}$} & \colhead{CMacIonize} & \colhead{Cloudy / MAPPINGS}}
\startdata
Dimensionality & 3D & 3D & 3D & 3D & 1D \\[2pt]
Grid & octree AMR + Cartesian DDA & uniform Cartesian & uniform Cartesian & Voronoi / adaptive Cartesian & spherical / plane-parallel \\[2pt]
Solution-driven refinement & yes (I-front) & no & no & static AMR & 1D shells \\[2pt]
Diffuse field & explicit MC default; OTS option & explicit MC & explicit MC & explicit MC & outward-only / OTS \\[2pt]
Recombination data & hydrogenic cascade (H$^{+}$, He$^{2+}$), Badnell 2023 (He$^{+}$, metals) + Milne $\alpha_1$ & Verner \& Ferland 1996 & MAPPINGS V & Verner / Badnell & Badnell / internal \\[2pt]
Metal ionization solve & ratio chain, one element at a time & ratio chain, one element at a time & MAPPINGS network & ratio chain, one element at a time & matrix, one element at a time (\textsc{Cloudy}) \\[2pt]
Metal feedbacks wired in & opacity, $\nelec$, photoheating, recombination and collisional-ionization cooling, free-free, line cooling & opacity, $\nelec$, photoheating, line cooling & MAPPINGS set & line cooling only & all, plus two-way charge exchange \\[2pt]
Dust / PAH & absorption + forward scattering + equilibrium + stochastic + PAH survival & absorption + emission & MAPPINGS grains & absorption + emission & detailed grains/PAH \\[2pt]
Parallelism & MPI-3 shared memory; static or on-demand packet schedule & MPI & MPI & task-based / MPI & OpenMP / serial \\[2pt]
Synthetic output & peel-off images, line maps, IR SED & line maps, images & line maps & line maps, images & 1D spectra \\
\enddata
\end{deluxetable*}

%==========================================================================
\section{Methods}
\label{sec:methods}

\MoCHII{} follows the conventions of its parent dust code: all input is a
single Fortran namelist (\code{par}); computation is in double precision;
parallelism is MPI-only with MPI-3 shared-memory windows for the grid;
output is HDF5 with a FITS fallback; and algorithm variants are selected
by procedure pointers at startup.  A run reads (or builds) a grid, persists
a gas leaf state, and iterates transport and equilibrium to convergence.

\subsection{The ionizing band and the source}
\label{sec:band}

The ionizing band spans photon energies $[E_{\min}, E_{\max}] =
[13.598, 100]$~eV.  In continuous mode, the source spectrum --- a Planck
function $B_\nu(T_\star)$ or a tabulated spectrum in shape or physical
(per-eV, per-Hz, per-\AA, per-$\mu$m) units --- defines an energy-luminosity
probability density.  A packet energy is drawn by inversion of its CDF and
the photoionization cross sections and excess energy are evaluated at that
drawn energy.  Thus the $n_\nu$ bins (typically 32) are a diagnostic tally
of $J_E$, not a transport discretization; changing their count cannot change
the physical solution.  Packets carry $L_{\mathrm{pkt}} =
L_{\mathrm{band}}/N$.

\paragraph{FUV extension and composite sources.}  With \code{par\%add\_fuv}
the band grows a second, lower segment of FUV bins on
$[E_{\mathrm{FUV}}, E_{\min}]$ below the unchanged ionizing bins.
Ground-state H and He cross sections vanish below their thresholds, so the
FUV photons interact only with dust, with the low-threshold metal stages,
producing the singly ionized carbon, sulfur, magnesium, silicon and iron
of the neutral zone (e.g.\ C\,\textsc{ii}, Mg\,\textsc{ii}) that the
ionizing-only band cannot, and with \HeI{} $2^3$S, whose photoionization
threshold is at 4.77~eV.  The extension follows the $2^3$S population
(Section~\ref{sec:heimeta}): whenever the run carries it --- by default,
and only when there is helium --- the extension is on unless the input
turns it off, and $E_{\mathrm{FUV}}$ defaults to 4.78~eV, the lower end of
the $2^3$S cross-section fit; otherwise the extension is off unless the
input turns it on, and $E_{\mathrm{FUV}}$ defaults to 6~eV.  The FUV photons are a
separate stratum of packets, by default a quarter as many as the ionizing
ones, each stratum carrying its own luminosity over its own count, so a
source whose FUV luminosity exceeds its ionizing one (14.6 times at
20~kK) does not take packets from the ionizing band and the ionizing-band
noise is that of a run without the extension.  The band is fed by any number of internal point sources
(0--16) plus, independently, one isotropic external radiation field.  The
external field is a band-integrated mean intensity $J$
[erg\,s$^{-1}$\,cm$^{-2}$\,sr$^{-1}$] entering through
the box faces or a bounding sphere; the inward flux is the hemisphere
integral of $J\cos\theta$, so the entering band power is $L = \pi J
A_{\mathrm{surf}}$, and entry directions are cosine-weighted (Lambert),
reproducing an isotropic interior field of energy density $u = 4\pi
J/c$.  Three analytic interstellar-radiation-field presets (Draine,
Habing, Mathis) are available for the FUV segment.

\subsection{Composable sources: a cloud in an external field}
\label{sec:external}

The source composition described above combines any number of internal point
sources with one external field, so a cloud can be modeled as it is actually
illuminated.  Figure~\ref{fig:external} is the limiting case of no internal
star: a uniform sphere ($\nH = 1.7\,\mathrm{cm^{-3}}$, $R = 1.8$~pc) bathed
in an isotropic external field with a $4\times10^4$~K blackbody shape,
ionized entirely from outside.  The resulting structure inverts that of a
central-source Str\"omgren sphere.  An ionized skin
(median $x_{\mathrm{HII}} = 0.92$ over the outer 0.3~pc) wraps a more neutral core
(median $\xHI = 0.94$ inside 0.5~pc), and the ionized-to-neutral crossover falls near
$r = 1.05$~pc.  The planar slice through the center and the shell-averaged
radial profiles both show the field working inward from the illuminated
surface.  A cloud embedded in an interstellar or circumgalactic radiation
field is therefore computed directly, without imposing the incident field as
a boundary condition on a separate one-dimensional model.

\begin{figure*}[t]
\centering
\includegraphics[width=0.95\textwidth]{figures/external_field.pdf}
% source: python/make_external.py, from examples/external_field
\caption{A uniform cloud ($\nH = 1.7\,\mathrm{cm^{-3}}$, $R = 1.8$~pc)
photoionized by an isotropic external radiation field entering the box faces
(a $4\times10^4$~K blackbody shape), with no internal source.  Left: a
planar slice of the neutral fraction $\xHI$ through the cloud center,
showing an ionized skin over a neutral core.  Right: the shell-averaged
radial profiles of $\xHI$ and $x_{\mathrm{HII}}$, with the field entering
at the surface $r = R$.  The ionization structure is the inverse of that
produced by a central source.}
\label{fig:external}
\end{figure*}

\subsection{Opacity and the zero-variance edge walk}
\label{sec:edge}

A packet of photon energy $E$ sees, in each leaf, the absorption coefficient
\begin{multline}
  \kappa(E) = \nH\big[\xHI\,\sigma_{\mathrm{HI}}(E) \\
   + y_{\mathrm{He}}\big(x_{\mathrm{HeI}}\sigma_{\mathrm{HeI}} +
     x_{\mathrm{HeII}}\sigma_{\mathrm{HeII}}\big)(E)\big] d_{\mathrm{cm}}
   \;(+\,\kappa_{\mathrm{dust}}(E)),
\end{multline}%
plus the metal photoionization opacity, with every cross section evaluated
at the packet's own energy rather than at a bin center.  Because the
ionizing band has no scattering (without dust) and the diffuse field, when
used, is emitted explicitly, each packet's path through a leaf is scored
\emph{analytically} along its ray --- the forced-first-flight estimator of
\citet{lucy1999}, made exact:
\begin{equation}
  \ell_{\mathrm{eff}} = w\,L_{\mathrm{pkt}}\,
  \frac{e^{-\tau_{\mathrm{in}}} - e^{-\tau_{\mathrm{out}}}}{\kappa(E)},
  \label{eq:edgewalk}
\end{equation}%
accumulated leaf by leaf until the domain edge.  No interaction sampling
is needed; the only Monte Carlo variance is in the direction and energy
sampling.  Each segment adds its $\ell_{\mathrm{eff}}$, weighted by the
coefficient of each rate at the packet's energy, to path estimators, which
after the reduction over ranks give for each absorber $i$ the
photoionization and photoheating rates
\begin{align}
  \Gamma_i &= \frac{1}{V d_{\mathrm{cm}}^2}\sum_{\mathrm{segments}}
  \ell_{\mathrm{eff}}\,\frac{\sigma_i(E)}{E},
  \nonumber\\
  \mathcal{H}_i &= \frac{1}{V d_{\mathrm{cm}}^2}\sum_{\mathrm{segments}}
  \ell_{\mathrm{eff}}\,\sigma_i(E)
  \left(1 - \frac{E_{\mathrm{th},i}}{E}\right),
  \label{eq:rates}
\end{align}
with $V$ the leaf volume in code units and, for a metal ion above its first
inner-shell edge, the heating coefficient summed over the absorbing
subshells (Section~\ref{sec:xsec}).  On request (\code{par\%jnu\_output},
or when the stochastic dust emission needs it) the same segments also fill
an $n_\nu$-bin mean-intensity array, $J_b$, which is written out as a
diagnostic spectrum and enters no rate.

\subsection{Grids: octree AMR and the Cartesian DDA backend}
\label{sec:grids}

\MoCHII{} inherits the \MoCafe{} octree.  The grid may be read as an AMR
leaf list, built from a namelist size ($n_x\times n_y\times n_z$) with an
optional uniform-density sphere or shell, or read from a 3D density cube.
Geometry is fully analytic behind accessor functions, so a leaf's center
and half-size are array reads for a general octree and raster-index decodes
for a uniform Cartesian (``car'') grid --- the same floating-point
expressions in both cases.  The car backend allocates no tree or neighbor
arrays at all (a saving that scales to $\sim$10~GB at $512^3$) and
traverses rays with a dedicated incremental Amanatides--Woo DDA walk, which
is the default for uniform grids.  A verification mode reproduces the
single-level octree traversal bit-for-bit, and the DDA agrees with it to
rounding, at up to twice the transport throughput.  The two backends share
all of the physics, transport, and peel-off code.

\subsection{Ionization and thermal equilibrium}
\label{sec:equil}

In each leaf, the ionization state is solved (\code{solve\_ion\_cell}) from the
ionization rates $R_i = \Gamma_i + \nelec C_i(T)$ and the case-A/B
recombination rates $\alpha_i$: hydrogen from
$\xHI R_{\mathrm{HI}} = x_{\mathrm{HII}}\alpha\nelec$ together with the
charge-exchange and on-the-spot terms, solved for
$u = \ln(n_{\mathrm{HII}}/n_{\mathrm{HI}})$ (Section~\ref{sec:compare}),
the He chain $n_{i+1}/n_i = R_i/(\alpha_{i+1}\nelec)$, closed
by $\nelec = \nH[x_{\mathrm{HII}} + y_{\mathrm{He}}(x_{\mathrm{HeII}} +
2x_{\mathrm{HeIII}})]$ and solved by a damped fixed point on $\nelec$.
$x_{\mathrm{HII}}$ and $x_{\mathrm{HeIII}}$ are stored beside $\xHI$ and
$x_{\mathrm{HeII}}$ rather than formed as differences from one, so a small
ionized fraction keeps its relative precision.  When
thermal balance is enabled, the temperature is updated once each iteration
by finding the root of the net rate $\mathcal{N}(T) = \mathcal{H}(T) -
\mathcal{C}(T)$ by bisection on $\log \Te$; the ionization state is
re-solved at every trial $T$.  The cooling $\mathcal{C}(T)$ assembles
recombination cooling, free-free (with the net ion charge $Z$, including
the trace metals at $Z_{\mathrm{ion}} = \mathrm{stage}-1$), collisional
ionization, \HI{} line cooling, and the trace-metal cooling sum
$\sum \nelec n_{\mathrm{ion}}\Lambda_{\mathrm{ion}}(T)$.  Optional secondary
ionization by fast photoelectrons redistributes the energy of every
\HI, \HeI{} and \HeII{} photoelectron above 30~eV between heating and
secondary \HI/\HeI{} ionization in partially neutral gas: the ionization
budget of \citet{shull1985} is carried to lower energies by the ratio of
mean energies per ion pair of \citet{dalgarno1999}, divided among the
neutral targets actually present, and the heat share is Dalgarno et al.'s
efficiency, with the ionized fraction taken from the current iterate of
the cell solve.  The photoelectron energy is accumulated during transport
on six energy nodes (30--1000~eV), so the partition, which depends on the
gas state, is applied exactly after the transport.

\subsection{The iteration loop and convergence}
\label{sec:iterate}

Each iteration zeroes the band tally, transports the stellar packets
(and, if enabled, rebuilds and transports the diffuse packets), reduces
the tally over MPI, forms the rate integrals, solves each leaf, and
refills the opacity from the new state; the stellar tally is recomputed
from zero each iteration.  The cell solves are divided among all ranks of
all nodes, each solving its own leaves from the reduced tallies; the solved
leaves are summed into every node's shared state, and the convergence
measures are reduced over all ranks, so every rank makes the same exit
decision.

Convergence is measured two ways every iteration, and
\code{par\%conv\_crit} selects which pair controls the loop: the largest
change found in any single cell ($\max|\Delta x_{\mathrm{HII}}|$,
$\max|\Delta \Te|/\Te$) and the volume-integrated measures
\begin{equation}
  \frac{\sum_{\mathrm{leaf}} V\,|\Delta x_{\mathrm{HII}}|}
       {\sum_{\mathrm{leaf}} V\,x_{\mathrm{HII}}},
  \qquad
  \frac{\sum_{\mathrm{leaf}} V\,\nelec\,|\Delta \Te|}
       {\sum_{\mathrm{leaf}} V\,\nelec\,\Te}.
\end{equation}
The cell criterion compares the hydrogen maximum with
\code{gas\_tol} and, when temperature is solved, the temperature maximum
with \code{gas\_tol\_te}; both tolerances default to $10^{-3}$.
It has no separate helium or metal-ion stopping test and needs no history.
The cell temperature denominator is the previous temperature; the volume
measures use the updated fractions, electron density, and temperature.
Ionization changes are measured after under-relaxation. The volume
criterion requires at least eight history entries for each tested ionization
measure, but tests temperature only on its current step, without extrapolation.
The default criterion is \code{cell}; \code{vol} selects the volume tests.
The volume measures address two stall modes of these single-cell maxima:
ionization-front jitter occupies little volume (the front shell is a few
thousand cells against $\sim$$10^5$ ionized ones, so its Monte Carlo
noise enters the $L_1$ norm suppressed by that ratio), and the $\nelec$
weight removes the vacuum and no-heating cells whose floor-pinned
oscillation dominates $\max|\Delta \Te|$ in FUV/PDR configurations.  On a
FUV dusty Str\"omgren case at $8\times10^6$ packets, where these single-cell
maxima stall at 0.1 and 6.4, the volume measures decay smoothly
to Monte Carlo floors of $2.7\times10^{-3}$ and $6.6\times10^{-3}$ by
iteration $\sim$27, and 30 further iterations leave the ionized volume
unchanged to 0.014\%.  The floors scale as $1/\sqrt{N_{\mathrm{packets}}}$,
and the volume criterion is the recommended production setting.

A test on the last change alone stops too early when the iteration
converges slowly, and in a radiation-bounded nebula it does: the fronts
advance by a bounded distance per iteration, and the volume measure decays
geometrically, $d_{n+1} \approx \rho\,d_n$, with $\rho$ close to one.  On
the HII40 benchmark of Section~\ref{sec:hii}, run for a fixed 80
iterations, $\rho \approx 0.86$ for $x_{\mathrm{HII}}$, so that six times
the last change is still to come, and the corresponding helium measure,
$\sum V|\Delta x_{\mathrm{HeI}}|/\sum V(1-x_{\mathrm{HeI}})$, is five times
larger than hydrogen's at the same iteration.  Stopping when the hydrogen
step fell below $5\times10^{-3}$ ended that run at iteration 37 with the
\HeI{} front 0.17~pc short of where it settles.  The volume criterion
therefore estimates the remaining change as
$\hat d_n\rho/(1-\rho)$ for hydrogen and helium, including He$^{2+}$
when its volume integral exceeds 1\% of that of ionized helium.
A least-squares fit $\ln d_k=a+bk$ over the last eight iterations gives
$\rho=\exp(b)$ and the fitted latest step $\hat d_n$.
When all steps in the window are positive and $\rho<0.95$, the estimated
remaining change must be below \code{gas\_tol}. Otherwise, the latest
measured step must be below that tolerance. An optional additional test
uses the signed counterpart of $d$,
\begin{equation}
  s_n = \frac{\sum_{\mathrm{leaf}} V\,(x_{\mathrm{HII},n} - x_{\mathrm{HII},n-1})}
             {\sum_{\mathrm{leaf}} V\,x_{\mathrm{HII},n}} :
\end{equation}
on the floor the leaves fluctuate about their fixed point and their changes
cancel in $s$, while a front that is still moving changes the ionized
volume with one sign every iteration. With
\code{par\%conv\_drift\_test = .true.}, the step test also requires
$|\overline{s}|/(\sigma_s/\sqrt{8})<2.365$, where $\sigma_s$ is the sample
standard deviation. This threshold is the two-sided 95\% Student value
for seven degrees of freedom, but correlated iteration histories do not
provide a formal 95\% convergence guarantee. The option defaults to
\code{.false.} as of 2026 October 18; the earlier calculations described
here used the additional test. It applies only to the volume ionization
tests when extrapolation is unavailable, not to temperature or the cell
criterion. Reading the floor off
the step ratios alone is not enough: a restart whose hydrogen front was
still advancing, with $\rho \approx 0.97$ and a step dominated by leaf
noise, passed after four iterations because one step ratio exceeded one;
continued, its $s$ gave $t = 8.8$--10, it took 40 further iterations to
satisfy the test above (the last 24 under that test), and its front moved
out by a further 0.012~pc.

The hydrogen balance residual is reported as a diagnostic and does not
control termination; neither diffuse luminosity nor metal-ion fractions
have separate outer stopping tests. The iteration limit is
\code{gas\_niter}; its default of zero performs transport without a gas
solve. Reaching a positive limit without convergence normally produces a
warning and output. Setting \code{require\_convergence = .true.} aborts
before final output instead.

After the loop, one further transport pass and gas solve are performed
when gas iteration is enabled. The written gas state is solved in the
written radiation field, which was transported through the preceding gas
state. The convergence flag retains the loop's decision; it is not retested
after this final solve. The output quantities \code{FINALDX} and
\code{FINALDTE} record the final solve's changes.

\paragraph{Parallel packet scheduling.}  The global packet indices may be
assigned statically or in dynamically dispatched batches.  Dynamic scheduling
reduces sensitivity to heterogeneous node contention because a rank that
finishes early requests additional packets.  The same scheduler is used for
stellar, diffuse-field, and peel-off transport.  Under pseudo-random
launching, the assignment determines the random-number stream used by each
rank; under the quasi-random launch of Section~\ref{sec:qmc}, packet
coordinates are fixed by the global index and are therefore independent of
the MPI task count and scheduling strategy.

\subsection{Explicit diffuse field}
\label{sec:diffuse}

When the explicit diffuse field is enabled, case-B OTS is replaced by
case-A rates plus explicit emission of ground-recombination packets in
three channels (\HII, \HeII, \HeIII) with rate $\alpha_1 = \alpha_A -
\alpha_B$.  Detailed balance fixes the emitted spectrum: per unit photon
energy the free--bound photon emissivity of each ground continuum is
\begin{equation}
  n(E) \;\propto\; E^{2}\,\sigma_{\mathrm{pi}}(E)\,
     \exp\!\left[-\frac{E - E_{\mathrm{th}}}{kT_e}\right],
  \qquad E \ge E_{\mathrm{th}},
  \label{eq:milne}
\end{equation}
evaluated with the \emph{same} ground-state photoionization cross section
the transport absorbs with, so a packet is absorbed with the cross section
it was emitted with.  Restoring the Milne prefactor turns the integral of
Equation~(\ref{eq:milne}) into $\alpha_1(T)$ itself, which recovers
$\alpha_{1s}(\HI, 10^4~\mathrm{K}) = 1.579 \times 10^{-13}$~cm$^3$\,s$^{-1}$
against the tabulated $1.58 \times 10^{-13}$.  \MoCHII{} takes $\alpha_1$
from that integral, so the case A/B split, the emitted spectrum and the
packet count all rest on one cross section; the code reports the integral
against the $\alpha_A - \alpha_B$ used for the emission rate without
rescaling either, and the two agree to $0.01$\% on all three channels.
Because packets
are equal-energy and their number follows the channel's \emph{energy}
luminosity, their energies are drawn from $E\,n(E)$ rather than $n(E)$,
which makes the represented photon rate exactly $n_e n_{\rm ion}\alpha_1 V$.
The cumulative distribution is tabulated in $x = (E-E_{\mathrm{th}})/kT$,
interpolated in $\log T$ and inverted once.  A near-threshold approximation
$E = E_{\mathrm{th}} + kT_e\,x$ with $x \sim \mathrm{Exp}(1)$, which is
Equation~(\ref{eq:milne}) with the $E^2\sigma_{\rm pi}$ weight dropped, is
retained as an option.  The emission table is rebuilt from the current state
each iteration; packets are transported by the same edge walk.  The
sampled photon energy is mapped to its bin by a binary search of the
stored bin-edge array, correct on both the aligned and the plain
logarithmic grids: an
energy exactly on a threshold edge falls in the bin \emph{above} it, the
physically correct side for a recombination photon emitted just above
threshold.

A fourth channel, on by default whenever the continuum is transported,
adds the \HeI{} excited-level recombination radiation.  Every case-B \HeII{} recombination reaches $1^1$S through
exactly one of the $n=2$ terms, so what the channel needs are the three
branch fractions; their absolute scale is already set by $\alpha_B(\HeII)$.
They are formed from one internally consistent set of effective
recombination coefficients,
$2.10\times10^{-13}\,T_4^{-0.381}$ ($2^3$S),
$2.06\times10^{-14}\,T_4^{-0.451}$ ($2^1$S) and
$4.17\times10^{-14}\,T_4^{-0.695}$ ($2^1$P) $\cmc$, from
\citet[][Table~1, with J.~S.~Mathis]{benjamin1999} as tabulated in K.~Wood's
\code{ionize} \citep{wood2004} and in \textsc{CMacIonize}
\citep{cmacionize2018}, normalized by their sum.  Taking all three from one
source matters more than the individual values, because a ratio of
coefficients with unrelated temperature dependences is not a branch fraction;
normalizing by the sum rather than closing one branch against an external
$\alpha_B$ keeps the mismatch between two independent data sets out of
whichever branch is solved last.  That sum is $0.91\,\alpha_B(\HeII)$ at
8000~K and $0.985$ at $10^4$~K, and the shortfall is reported at start-up
rather than absorbed.  At 8000~K the fractions are $0.762 : 0.076 : 0.162$
against the $0.750 : 0.083 : 0.167$ of the low-density statistical weights
($0.741 : 0.076 : 0.183$ at 5000~K, $0.771 : 0.076 : 0.153$ at $10^4$~K): a
modest temperature-dependent correction, not a reordering.  The $2^1$S
photons are drawn from the tabulated \citet{drake1969} two-photon
distribution rather than flat above the \HI{} edge; that shape also fixes the
H-ionizing fraction, 0.5564 photons per decay with a mean in-band energy of
16.11~eV.  Together the two replace the 19.200~eV in 0.9634 photons that the
statistical weights and a flat two-photon spectrum emit per case-B \HeII{}
recombination by 19.223~eV in 0.9663 photons at 8000~K.  The 584~\AA{}
photon of the $2^1$P branch is a resonance photon, and two irreversible
outcomes compete during its trapped random walk: absorption by a hydrogen
atom, which photoionizes it, or the 2.06~$\mu$m decay
$2^1$P\,$\to$\,$2^1$S of the helium atom that last absorbed it (about one
scattering in a thousand ends that way), after which the decay leaves as the
two-photon continuum instead.  \MoCHII{} splits the
branch with the escape probability that \citet{wood2004} and
\textsc{CMacIonize} \citep{cmacionize2018} adopt,
$p_{584} = [1 + 77\,(y_{\rm He}/0.1)\,x(\mathrm{He^{0}})/(\Te^{1/2}\,x(\mathrm{H^{0}}))]^{-1}$
with $\Te$ in K and $y_{\rm He} = n_{\rm He}/n_{\rm H}$ (their 77 assumes
He/H $=0.1$, their eq.~17); it is a property of the emitting leaf, so it folds into the
branch weights of the emission table and costs no further random variate.  A
converted decay radiates 20.62~eV as the two-photon pair, the remaining
0.60~eV leaving as the non-ionizing 2.06~$\mu$m $2^1$P\,$\to$\,$2^1$S line.
The \HeI{} two-photon nebular continuum is written from the same exits,
$r_{2q} = \nelec\, n(\mathrm{He^+})\,\alpha_B(\HeII)\,[f_{1S} +
(1-p_{584}) f_{1P}]$ per unit volume (plus the collisional transfers from
$2^3$S when that population is coupled), so the transported packets, the
on-the-spot sources and the tabulated continuum describe one nebula.
In both benchmarks most of the branch converts, so the channel loses a
fraction of its photons and softens; at HII20 the softer spectrum leaves
little \HeII{} to recombine, and the channel carries a few percent of the
diffuse field rather than the $\sim$20\% it carries at HII40.  With the
$2^3$S population coupled (Section~\ref{sec:heimeta}) the triplet exit is
no longer a fixed fraction: of the atoms that reach $2^3$S, those that
decay radiate the 19.82~eV photon, those transferred by electron impact to
$2^1$S or $2^1$P leave through the singlet exits, and those photoionized,
collisionally ionized or Penning-ionized emit nothing, so the exits depend
on the density and on the local field.  A fifth channel carries the
He$^{2+}$ case-B cascade --- the free--bound continuum into $n \ge 2$, whose
threshold at 13.604~eV lies just above the hydrogen edge, the 40.8~eV
\HeII{} Ly$\alpha$ and the $2^2$S two-photon continuum --- all of which
ionize hydrogen, with the branching taken from the level-resolved cascade of
Section~\ref{sec:recomb} (with \code{par\%lyman\_lines = 'escaping'} every
\HeII{} $n\mathrm{p}\to1\mathrm{s}$ line leaves at its own energy, and the
band must then reach 54.36~eV); it is on with the transported continuum
whenever the band reaches 40.8~eV, and it matters only where He$^{2+}$ exists, which the
benchmarks of Section~\ref{sec:hii} scarcely have.

In the on-the-spot mode none of these photons is transported, and
each is absorbed in the leaf that made it.  The \HeI{} and \HeII{}
recombination photons are then shared as in \citet[][Eqs.~2.24--2.25 and
\S2.5]{osterbrock2006}: a ground-capture photon goes to each absorber in
proportion to $n_X\sigma_X$, averaged over its Milne spectrum
(Equation~\ref{eq:milne}) by a four-point Gauss--Laguerre rule, and every
excited-cascade photon below the \HeI{} edge goes to hydrogen, with the same
branch fractions and exits as the transported channels.  The \HeI{}
cascade gives hydrogen 0.966 ionizing photons per recombination when every
584~\AA\ photon ionizes it, the limit in which \citet{osterbrock2006} find
$\simeq0.96$.  Each absorbed photon heats the gas by its energy above the
absorber's threshold, less, for a ground capture, the kinetic energy the
capture removed.  Before this treatment the on-the-spot mode gave every
\HeI{} recombination photon back to helium; on the Lexington HII40 model
the change raises $\langle T[N_{\mathrm{p}}N_{\mathrm{e}}]\rangle$ by
178~K and $L(\mathrm{H}\beta)$ by $7.8\%$ and moves the hydrogen front out
by $0.16$~pc, to within $0.02$~pc of the run with the continuum transported.
The approximation has a limit of its own: a cell can absorb on the spot only
as many of these photons as the recombinations that make H$^0$, and in a
fully doubly ionized zone the He\,\textsc{ii} cascade outnumbers them once
He/H exceeds a value between $0.12$ and $0.15$ (it reaches $0.92$ of them at He/H $=0.1$).
Beyond that no local balance exists; \MoCHII{} then keeps $x(\mathrm{H^0})
\approx 0$, heats with the photons the supply can absorb, reports the leaf
through the hydrogen balance residual, and such nebulae require the
transported continuum.

At HII40 the \HeI{} front sits $0.033$~pc inside the \textsc{Cloudy}
c25.00 one, $4.272$ against $4.305$~pc, or $0.8\%$ in radius (both from
volume-weighted 0.01~pc shells).  An earlier
version of this comparison found $0.18$~pc and bounded it candidate by
candidate --- the diffuse emission and its transport, the \HeI{}
photoionization and recombination data, the metal opacity, the photon
budget --- without finding one large enough.  The difference was the
stopping rule of Section~\ref{sec:iterate}: the hydrogen step had fallen
below its tolerance while the ionization structure was still converging, and
run to the drift test the front moves outward by $0.17$~pc to the position
above.  The sensitivities measured in that exercise were differences
between runs stopped by the earlier rule, and are not quoted here.

It did expose one defect, which stands.  $\alpha_1$ had been taken from a
level-resolved fit \citep{mao2016} while $\alpha_A$ came from the Badnell
total, so $\alpha_B = \alpha_A - \alpha_1$ crossed two unrelated
determinations; \textsc{Cloudy} c25.00 \citep{cloudy2025}, which derives the
same quantity level by level from its own cross sections through the Milne
relation, puts that fit $2.4\%$ low for \HeI.  Replacing it by the Milne
integral over the transport cross sections brings all three channels within
$0.35\%$ of \textsc{Cloudy} and the start-up cross-check from $3.4\%$ to
$0.01\%$.

The helium recombination data have also been compared network to network.
At $10^{4}$~K,
$\alpha_A$(\HeI) is $4.358 \times 10^{-13}\,\cmc$ in \textsc{Cloudy}'s
explicit sum over levels, $4.372$ in \MoCHII, $4.263$ in the effective
coefficients of \citet{benjamin1999} used by \citet{wood2004} and
\textsc{CMacIonize}, and $4.839$ in \textsc{Mocassin}'s Verner \& Ferland
(1996) fit.  No code truncates the cascade --- all three treatments run over every
level --- and the $11\%$ spread is confined to \HeI: for the one-electron
species, whose case-B rate \MoCHII{} takes from a level-resolved cascade
(Section~\ref{sec:recomb}), $\alpha_A$ agrees with \textsc{Cloudy}'s level
sum to $0.04\%$ (\HI) and $0.02\%$ (\HeII), where an exact analytic solution
leaves nothing to disagree about, so what spreads is the two-electron atomic
data itself.  \MoCHII{} sits within $0.3\%$ of \textsc{Cloudy}, which also
means any helium comparison against \textsc{Mocassin} inherits that
difference.  The \HeI{} excited levels a resolved model atom carries lie
below \MoCHII's 13.598~eV band floor and are worth $0.2$--$0.6\%$ of extra
ionization.  At HII40 the diffuse-band luminosity converges to
$\sim$$0.5\,L_\star$.

\subsection{Quasi-random photon launching}
\label{sec:qmc}

By default, \code{par\%launch\_sequence = 'sobol'} draws the launch
variables from scrambled Sobol sequences \citep{joekuo2008,burley2020};
\code{'random'} selects the Mersenne Twister. Stellar variables include
source component, energy, direction, and position as required by the source
geometry. Diffuse emission uses a separate scramble stream with nine fixed
coordinates: emitting leaf, three positions within that leaf, two direction
coordinates, emission channel, and two branch or energy coordinates.
The point is addressed by the global packet number within its stream,
independently of MPI scheduling. The scramble is fixed across gas iterations;
changing \code{qmc\_seed} provides another realization for error assessment.
Dust-scattering events and Russian roulette retain pseudorandom sampling.

Without dust scattering, the path contribution is integrated analytically
conditional on the launch variables (Eq.~\ref{eq:edgewalk}), so improving
launch sampling can reduce the radiation-field error. Cell boundaries,
channel selection, and ionization fronts can nevertheless limit the gain.
Diffuse packet counts follow the luminosity ratio and generally are not
powers of two: arbitrary Sobol prefixes are usable, but do not retain the
full digital-net balance guarantee of a $2^m$ sample. Reusing a launch set
can also reduce iteration changes without establishing accuracy of the
final solution.

Figure~\ref{fig:qmc} retains the historical fixed-state stellar
analytic-attenuation test, with dust scattering and transported recombination
emission disabled. The archived analysis quoted mean-squared-error ratios
of $28.8$, $66.7$, $143$, and $231$ at $N=2^{17},2^{18},2^{19},2^{20}$,
respectively. These are ratios of squared errors relative to the reference,
not estimates of sampling variance with bias removed. The plotted
$N^{-1/2}$ and $N^{-1}$ lines are reference slopes for that experiment,
not general convergence guarantees for diffuse radiation.

These numerical results are historical and remain unverified for the current
continuous-energy implementation: the figure script evaluates its reference
at energy-bin centers. A new comparison must use the same continuous source
spectrum and cross sections as transport before these gains, the earlier
bias estimates, or equal-time performance claims can characterize the current
code. Isolating the diffuse benefit additionally requires holding stellar
sampling fixed while varying diffuse sampling across independent realizations.

\begin{figure}[t]
\centering
\includegraphics[width=\columnwidth]{figures/qmc_variance.pdf}
% source: python/make_qmc.py, from tests/qmc/v3/
\caption{Quasi-random photon launching.  Cell-wise root-mean-square error
in $\Gamma_{\mathrm{HI}}$ on the fixed-state analytic-attenuation test, as a
function of the launched photon count $N$, for the ordinary Monte Carlo
(\texttt{random}) and Owen-scrambled Sobol (RQMC) launches, with the
$N^{-1/2}$ and $N^{-1}$ reference slopes.  The boxed labels give the
ratio of mean squared errors relative to the historical reference at each
$N$. These archived results require a continuous-energy reference before
being used to assess the current implementation; they do not measure
diffuse-field convergence.}
\label{fig:qmc}
\end{figure}

\subsection{Peel-off imaging}
\label{sec:peel}

When observer peeling is enabled, the final transport pass uses it
(the \MoCafe{} convention adapted to the band), including when an
unconverged run is allowed to produce output.  Every emitted packet
peels a direct contribution $L_{\mathrm{pkt}}\,w\,e^{-\tau}/(4\pi r^2)$
toward each observer, and, when dust scattering is sampled, every
interaction peels a scattered contribution with the Henyey--Greenstein
phase function.  The optical depth to the box edge is integrated by a
tally-free walk, so the images carry exactly the transport's extinction.
Images are TAN projections written as surface brightness, split into
direct/scattered and ionizing/FUV channels; with a spectral cube option
the channel axis becomes the band bin for hardness maps.  Leaf-resolved
line and continuum emissivities are also written, from which
flux-conserving line and field maps are produced.

Figure~\ref{fig:peel} is an example: the FUV channel of a dusty,
radiation-bounded \HII{} region with a central star, seen by the observer.
The direct image is the attenuated stellar point source (peak
$\log_{10} I = 1.95$ in $\mathrm{erg\,s^{-1}\,cm^{-2}\,sr^{-1}}$), and the
scattered image is the dust halo it illuminates (peak $0.24$, covering some
$27\%$ of the pixels).  The halo is the observer-dependent component that a
one-dimensional or angle-averaged calculation does not provide, and the sum
of the two is the surface brightness an observer measures.  The ionizing band contributes almost nothing to the
scattered image here, because in a bounded nebula it is absorbed before it
can scatter out.  Alongside the extincted image the code also writes the
bin-resolved image cube and an unattenuated (\code{direc0}) image, the
latter giving a line-of-sight optical-depth map as its ratio to the
extincted one.

\begin{figure*}[t]
\centering
\includegraphics[width=\textwidth]{figures/peel.pdf}
% source: python/make_peel.py, from tests/peel/peel_scat
\caption{Peel-off synthetic imaging of a dusty, FUV-illuminated \HII{}
region with a central star, in the FUV channel.  Left: the direct stellar
image, an attenuated point source.  Center: the dust-scattered halo, the
observer-oriented contribution.  Right: their sum, the synthetic
observation.  All three panels share one logarithmic scale of the specific
intensity $I$ integrated over the channel's band, so the extended scattered
halo is directly comparable to the compact direct source.}
\label{fig:peel}
\end{figure*}

%==========================================================================
\section{Atomic data}
\label{sec:atomic}

The design principle for the atomic data is that adding an ion is a
\emph{data operation, not a code operation}.  All rate coefficients are
fitted offline by Python programs from CHIANTI v11.0.2
\citep{dere1997,delzanna_chianti} and the Strathclyde/Badnell tables, and
are evaluated online by fixed analytic formulae.  Every coefficient file
carries a provenance header (source, version, fit range, maximum fit
error), so the coefficients are auditable and the Fortran only evaluates
formulae.

\subsection{Photoionization cross sections}
\label{sec:xsec}

The one-electron ions \HI{} ($Z=1$) and \HeII{} ($Z=2$) use the exact
ground-state hydrogenic cross section \citep{osterbrock2006}
\begin{gather}
  \varepsilon = \sqrt{E/E_{\mathrm{th}}-1},
  \nonumber\\
  \sigma(E) = \frac{\sigma_0}{Z^2}\left(\frac{E_{\mathrm{th}}}{E}\right)^{4}
  \frac{\exp\!\big(4 - 4\arctan\varepsilon/\varepsilon\big)}
       {1 - e^{-2\pi/\varepsilon}},
  \label{eq:sighyd}
\end{gather}%
with $\sigma_0 = 6.30\times10^{-18}\,\mathrm{cm^2}$ and $E_{\mathrm{th}}$
the true ionization potential (13.598 and 54.416~eV), tied to the physical
edge so it stays aligned with the band grid.  \HeI{} uses the single-shell
analytic fit of \citet{verner1996}:
\begin{gather}
  x = \frac{E}{E_0} - y_0,\qquad z = \sqrt{x^2 + y_1^2},
  \nonumber\\
  \sigma(E) = \sigma_0\left[(x-1)^2 + y_w^2\right]
  z^{-Q}\left(1 + \sqrt{z/y_a}\right)^{-P},
  \nonumber\\
  Q = 5.5 - \tfrac{1}{2}P,
  \label{eq:vfky}
\end{gather}%
with parameters parsed directly from the authors' \code{phfit2.f} table.

For the metal ions \MoCHII{} evaluates what \code{phfit2.f} itself
evaluates, which is not one fit at all energies.  Below the first
inner-shell threshold $E_{\mathrm{inn}}$ of the ion the outer-shell fit of
Eq.~(\ref{eq:vfky}) is used; at and above it the cross section is the sum of
the \citet{verner1995} fits of every occupied subshell $s$ whose threshold
has been passed,
\begin{gather}
  y = E/E_{0,s},
  \nonumber\\
  \sigma(E) = \sum_{s:\,E \ge E_{\mathrm{th},s}} \sigma_{0,s}
  \left[(y-1)^2 + y_{w,s}^2\right] y^{-(5.5 + l_s - P_s/2)}
  \nonumber\\
  \times\left(1 + \sqrt{y/y_{a,s}}\right)^{-P_s},
  \label{eq:vy95}
\end{gather}%
with $l_s$ the orbital quantum number of the subshell.  \code{phfit2} sets
$E_{\mathrm{inn}} = 0$ for $Z = 15$, 17, 19 and for $Z > 20$ except iron,
the elements absent from the outer-shell table (Opacity-Project gaps), so
for chlorine the subshell sum is used at every energy.  Seven registry ions
have their first inner-shell edge inside the transported band:
Mg\,\textsc{i} and Mg\,\textsc{ii} ($2p$; 54.9 and 65.7~eV),
Fe\,\textsc{i}--\textsc{iii} ($3p$; 66.0, 76.2 and 87.1~eV), and
Ca\,\textsc{i} and Ca\,\textsc{ii} ($3p$; 34.4 and 40.9~eV).  Above those
edges the outer-shell fit alone, which earlier versions of the code used up
to 100~eV, leaves out the inner subshell entirely: at 85~eV it gives
$0.050$~Mb for Mg\,\textsc{ii} where the subshell sum gives $6.68$~Mb, and
at 50~eV $0.090$~Mb for Ca\,\textsc{i} where the sum gives $1.43$~Mb.  Each
subshell releases its photoelectron with $E - E_{\mathrm{th},s}$, so the
photoheating coefficient is $\sum_s \sigma_s (1 - E_{\mathrm{th},s}/E)$ over
the absorbing subshells.  The element files carry the outer-shell fit,
$E_{\mathrm{inn}}$ and the subshell rows exactly as parsed from
\code{phfit2.f}, and a standing test compares the cross section and the
heating coefficient of every registry ion with \code{phfit2.f} itself at
9112 energies between 5 and 1000~eV, to a worst relative difference of
$2.3\times10^{-6}$.  The exoplanet-atmosphere code EXHALE evaluates the
same \code{phfit2} prescription.

What the code does not yet follow is the fate of the inner-shell vacancy.
A $2p$ hole in Mg\,\textsc{i}, and a $3p$ hole in Ca\,\textsc{i} or in
Fe\,\textsc{i}--\textsc{iii}, lies above the threshold for a second
ionization and autoionizes: the \citet{kaastra1993} electron yields, as
tabulated in \textsc{Cloudy} c25.00, give a second electron with
probability 0.995--0.9999 for these five vacancies, while the
Mg\,\textsc{ii} $2p$ and Ca\,\textsc{ii} $3p$ vacancies give one.
\MoCHII{} raises the ion by one stage and assumes that the excitation
energy $E_{\mathrm{th},s} - E_{\mathrm{th}}$ leaves the gas as fluorescence.
The nearest-neighbor ratio chain (Section~\ref{sec:compare}) cannot
represent a two-stage jump, so for these five ions the stage fractions and
the heating above the inner edges are approximate, while the cross section,
and hence the opacity, is not.  Mg, Ca and Fe are trace species in an
\HII{} region and are absent from the Lexington abundance set of
Section~\ref{sec:hii}; the size of the error has not been measured, and
removing it requires the Auger stage matrix planned for harder spectra.

\subsection{Recombination}
\label{sec:recomb}

The default recombination model takes, for all three ions, the ground-level
($n=1$) direct-capture rate $\alpha_1$ from the Milne relation applied to
the same photoionization cross sections the transport absorbs with
(Section~\ref{sec:diffuse}), evaluated at run time from the
seven-coefficient form of \citet{mao2016} with coefficients refitted to
that integral, and the rest from the best determination for each ion.  For
the one-electron ions H$^{+}$ and He$^{2+}$ the case-B rate $\alpha_B$ is
the zero-density radiative cascade resolved in $n$ and $l$ ($n \le 499$
plus a power-law tail), computed with the hydrogenic routines of
\citet{storey1991}, tabulated over $\log T = 1.5$--7.0 and interpolated by
a cubic in $\log$--$\log$ with a continuous derivative, since the
recombination cooling below differentiates it; it lies within $0.1\%$ of
the total of \citet{storey1995} at $10^4$~K for both ions.  Case A adds the
ground term, $\alpha_A = \alpha_B + \alpha_1$.  For $\HeII\to\HeI$ the
total $\alpha_A$ is the \citet{badnell2006} radiative recombination fit,
\begin{gather}
  \alpha_{\mathrm{RR}}(T) =
  \frac{A}{\sqrt{T/T_0}\big(1+\sqrt{T/T_0}\big)^{1-b}
           \big(1+\sqrt{T/T_1}\big)^{1+b}},
  \nonumber\\
  b = B + C\,e^{-T_2/T},
  \label{eq:badnellrr}
\end{gather}%
plus the Badnell three-term dielectronic sum, and case B removes the ground
term, $\alpha_B = \alpha_A - \alpha_1$.  The hydrogenic ions were formed
that way too until the cascade replaced it, and the subtraction amplified
the $0.6\%$ by which the fitted H$^{+}$ total is high at $10^4$~K into a
$1.1\%$ excess in $\alpha_B$: a fit to a total, less an exact part,
carries the fit's absolute error on a smaller quantity.  Either way the
construction mirrors how \textsc{Cloudy} splits case A and case B (a
level-resolved sum and its ground term) and reuses the same modern data.
An alternative setting selects the \citet{huignedin1997} fits; in case B
the two agree to within $1\%$ for \HII{} and \HeIII{} over 5000--20000~K,
the one substantive difference being $\HeII\to\HeI$, where the Badnell
total ($\alpha_A = 4.37\times10^{-13}\,\cmc$ at $10^4$~K) coincides with
\textsc{Cloudy}, the older \citet{huignedin1997} case-B value is $5\%$
lower than \MoCHII's, and the \textsc{Mocassin} Verner \& Ferland
$\alpha_A$ ($4.84$) is the high outlier.
For the metal cascade, total recombination is the \citet{badnell2006}
radiative recombination plus the Badnell-group dielectronic fits
$\alpha_{\mathrm{DR}}(T) = T^{-3/2}\sum_i c_i e^{-E_i/T}$
\citep{badnell2003}, read from the Strathclyde tables (verified md5-identical
to the copies \textsc{Cloudy} redistributes), with a CHIANTI fallback for
the few ions Badnell does not tabulate.

The recombination cooling is formed from the same coefficients.  For a
Maxwellian electron distribution and any capture cross section ---
radiative or dielectronic, summed over all levels --- differentiating
$\alpha = \langle\sigma v\rangle$ with respect to $T$ gives exactly the
kinetic energy that the captured electrons remove,
\begin{equation}
  \beta(T) = kT\left[\tfrac{3}{2}\,\alpha(T) + \frac{d\alpha}{d\ln T}\right],
  \label{eq:betarec}
\end{equation}%
so \MoCHII{} evaluates $\beta$ for \HII, \HeII{} and \HeIII{} from the
case-A or case-B coefficient that the ionization balance uses, with the
derivative taken by a centered difference in $\ln T$, and the balance and
the cooling cannot refer to different coefficients.  The
\citet{huignedin1997} cooling fits used before were fitted to a different
recombination coefficient and lie $0.8\%$ below Eq.~(\ref{eq:betarec})
evaluated on \MoCHII's own case-B \HII{} coefficient at $10^4$~K; the
customary $kT\alpha$ for \HeII{} is the mean energy removed only when
$\alpha \propto T^{-1/2}$, and on the Badnell case-B coefficient each \HeII{}
recombination removes $0.70\,kT$ at $10^4$~K, so $kT\alpha$ overstated that
channel $1.43$ times.  The same relation applied to each metal transition,
with the Badnell total, gives the kinetic energy lost to metal
recombination, and collisional ionization of a metal stage removes its
ionization energy; both losses enter the thermal balance together with the
metal photoheating, as the two sides of one energy budget
(\code{par\%metal\_heat}).

\subsection{Collisional ionization and charge exchange}
\label{sec:ci}

Collisional ionization uses the \citet{voronov1997} fit
\begin{equation}
  k(T) = A\,\frac{1 + P\sqrt{U}}{X + U}\,U^{K} e^{-U},
  \qquad U = \frac{\Delta E}{kT},
\end{equation}%
with an optional hybrid that multiplies each rate by the constant
\citet{dere2007}/Voronov ratio of that element and stage; the difference matters
only in shielded, low-ionization zones where collisions are the sole
metal-ionization channel.  Charge exchange with hydrogen uses
\citet{huang2023} Table~4 for the first stages except nitrogen and sulfur,
and the
\citet{kingdon1996} fits $k = a\,T_4^{b}(1 + c\,e^{d\,T_4})$ for the higher
stages (recombination direction), except for $\mathrm{O}^{2+}$ and
$\mathrm{N}^{2+}$ below.  The tabulated higher-stage rates are
indexed by the \emph{product} ion of the reaction
$\mathrm{X}^{i+} + \mathrm{H}^{0} \to \mathrm{X}^{(i-1)+}$, a convention worth
stating explicitly because it is easy to misread by one ionization stage and
because the doubly and triply ionized metal fractions are sensitive to it;
the indexing used here is checked against \textsc{Mocassin}'s ionization loop
and \textsc{CMacIonize}.  The same product-versus-reactant hazard appears in
the first-stage table itself.  The two oxygen rows of \citet{huang2023}
Table~4 are printed with their row labels interchanged: measured against the
detailed-balance value
$[g(\mathrm{X}^{+})g(\mathrm{H}^{0})/g(\mathrm{X}^{0})g(\mathrm{H}^{+})]\,
e^{-\Delta E/kT}$, the ratio of the two directions is $1.47$ at $8000$~K for
the printed rows and $0.910$ once they are exchanged.  The corrected
assignment agrees to $1$--$3\%$ with the \textsc{Cloudy} c25.00
charge-transfer fit of \citet{stancil1999}, the determination that the oxygen
footnote of \citet{huang2023} itself cites, so the data were right and only
the labels were transposed.  A detailed-balance test of every transition-1
pair over $3000$--$20000$~K is now a standing gate.  Its first version
exempted pairs whose smaller rate was below $10^{-15}\,\cmc$, and that
exemption hid a second defect: the nitrogen rows of \citet{huang2023}
Table~4 make the endothermic $\mathrm{N} + \mathrm{H}^{+}$ fall with
temperature, to $8.9\times10^{-20}\,\cmc$ at $10^4$~K, while
$\mathrm{N}^{+} + \mathrm{H}$ sits $70$--$86$ times below the
\citet{kingdon1996} value.  Nitrogen now takes \citet{kingdon1996}
(their Tables~1 and 3, one calculation and its detailed-balance reverse;
$1.54\times10^{-12}$ and $1.01\times10^{-12}\,\cmc$ at $10^4$~K), and the
gate has no rate floor.

Sulfur's first stage is resolved by the product term of S$^+$: \citet{zhao2005}
give the forward rate into $^4$S, $^2$D and $^2$P separately, and the reverse
is each channel's detailed balance weighted by the term's actual population
from the [S\,\textsc{ii}] $n$-level atom.  A reverse formed from the summed
forward rate by one detailed balance, as \citet{huang2023} do for the rates
they give only in the exothermic direction, presumes the excited terms in
LTE; at $10^4$~K and $n_e\rightarrow0$ that overstates the reverse from the
ground term $1.6{\times}10^4$-fold ($3.7{\times}10^{-13}$ against
$2.3{\times}10^{-17}\,\cmc$).  The Mg, Si and Fe first-stage reverse rates
are of that form and have no term-resolved forward rate to replace them;
they are upper bounds.

For $\mathrm{O}^{2+} + \mathrm{H}$ and $\mathrm{N}^{2+} + \mathrm{H}$ the
rate coefficients are those computed by \citet{barragan2006} (their
Table~1, $k_1$, and Table~2, $k_4$), the calculation that \textsc{Cloudy}
and CHIANTI also use, interpolated log--log in $T$ over the tabulated
$10^2$--$10^5$~K and held at the end values outside it.  The
\citet{kingdon1996} fits they replace are $1.22$ and $1.56$ times higher for
$\mathrm{O}^{2+}$ at $10^4$ and 5000~K, and $0.57$--$0.64$ times the
tabulated rate for $\mathrm{N}^{2+}$ over the same range.

Helium exchanges charge with hydrogen as well
(\code{par\%helium\_charge\_exchange}, on by default).
$\mathrm{He}^{+} + \mathrm{H} \to \mathrm{He} + \mathrm{H}^{+}$ proceeds
through two channels that the calculation of \citet{zygelman1989} keeps
apart.  The radiative channel is
$1.20\times10^{-15}\,(T/300)^{0.25}\,\cmc$
\citep[][Table~1, reaction 19]{stancil1998}, one quarter of the ``Direct''
column of Zygelman et al.'s Table~I, the quarter being the probability of
approach in the singlet state \citep[][Eqs.~4--6]{courtney2021}.  Over
$10^3$--$10^4$~K it agrees to 2\% with the recalculation of
\citet{courtney2021}; above $1.5\times10^4$~K the $T^{0.25}$ extrapolation
lies 7--13\% above it.  The
non-radiative channel, the ground-to-ground exchange, is taken from the
fully quantal calculation of the reverse reaction by \citet{loreau2014}
(their Eq.~3 and Table~2) through microscopic reversibility, with the
statistical-weight ratio $1/4$ supplying the singlet approach.  Between
the reverse threshold and the first quantal energy (15~eV/u) the cross
section is taken to vanish as the square root of the energy above
threshold, i.e.\ $\sigma_f v$ constant at low forward energy, as for
capture on the ion--induced-dipole potential with an energy-independent
transfer probability; the fit carried through the threshold and a zero
cross section there bracket this form.  The rate is then nearly constant,
$(0.95$--$1.02)\times10^{-16}\,\cmc$ over $10^3$--$3\times10^4$~K, 3\% of
the radiative channel at $10^4$~K, where it agrees with the non-radiative
cross sections of \citet{zygelman1989} Maxwell-averaged and multiplied by
$1/4$ ($0.93\times10^{-16}$).  The \citet{kingdon1996} fit used before is
$93\times$ larger there: at $5\times10^4$--$10^5$~K it reproduces
Zygelman's singlet-conditional rate without the $1/4$, and below
$2\times10^4$~K it exceeds the bound that Zygelman's own cross sections set
by at least a factor of 11, so no constant rescaling mends it.  Above
$\sim$$3\times10^4$~K the two quantal calculations disagree by up to $10^3$
in the cross section, and the rate there is uncertain by that much.  $\mathrm{He}^{2+} + \mathrm{H} \to \mathrm{He}^{+} +
\mathrm{H}^{+}$ takes the constant $10^{-14}\,\cmc$ of \citet{kingdon1996}.
The two reverse reactions follow by detailed balance --- for
$\mathrm{He} + \mathrm{H}^{+} \to \mathrm{He}^{+} + \mathrm{H}$ of the
non-radiative channel alone, since the radiative channel has no collisional
reverse.  The direct calculation of that endothermic direction by
\citet{kimura1993} is not used: its charge-transfer cross sections start
at 20~eV, above the 10.99~eV threshold, and its tabulated rates imply an
apparent activation energy of 4.3~eV between 8000 and $10^4$~K, well below
what the threshold permits, so they are not reliable at nebular
temperatures.  Every reaction enters both
the hydrogen and the helium balance, so the charge it moves is conserved in the pair.
The energy a reaction moves between binding energies enters the thermal
balance as the kinetic energy of its products,
$\Delta E_{\rm kin} = [{\rm IP}(\text{acceptor's product stage}) -
{\rm IP}(\text{donor})] - E_{\rm rad}$, with $E_{\rm rad}$ the
excitation of the product state that leaves as radiation (for the radiative
He$^{+}$ + H channel, the emitted photon); an endothermic reaction takes
$|\Delta E_{\rm kin}|$ from the gas, the sign convention of
\citet[][\S2]{kingdon1996}.  The non-radiative He$^{+}$ + H channel thus
releases 10.99~eV per event.  Product states are identified from CHIANTI
level energies where a defect is printed for the fit used, and are stated
for every carried reaction in the element files.  One reaction needs a
temperature-dependent defect: Mg + H$^{+}$ populates both Mg$^{+}$ 3p
($+1.52$~eV) and the endothermic Mg$^{+}$ 4s ($-2.70$~eV) in the
proportion of the two channel rates of \citet{allan1988}, so it takes
$1.26$~eV per event from the gas at $10^{4}$~K.  On the HII40 benchmark the
heating raises $\langle T[N_{\mathrm{p}}N_{\mathrm{e}}]\rangle$ by 1.4~K.
On the two Lexington benchmarks of Section~\ref{sec:hii}
the helium--hydrogen exchange changes $\langle T[N_{\mathrm{p}}N_{\mathrm{e}}]\rangle$
by $0.02$~K or less and $L(\mathrm{H}\beta)$ by $2\times10^{-6}$ or less,
because He$^{+}$ and H$^{0}$ scarcely coexist there.

\subsection{Cooling and $n$-level diagnostic data}
\label{sec:cooldata}

Collisional line cooling enters the code at two levels, both built from the
same CHIANTI v11.0.2 model of each ion (level energies from the observed
term structure, radiative $A$-values, and effective collision strengths
descaled after \citealt{burgess1992}).  The thermal loop needs a cooling
rate at every trial temperature of the bisection, and therefore uses a
closed-form fit suppressed to the local density (below); the diagnostic
output needs density-dependent emissivities, and therefore solves the level
populations.  Both are the same statistical-equilibrium solution of each
ion, so the thermal balance and the emergent line list are one physics.

The thermal-loop coefficient is a multi-exponential fit
\begin{equation}
  \Lambda_{\mathrm{ion}}(T) = T^{-1/2}\sum_{i=1}^{N} A_i\,e^{-T_i/T}
  \quad [\mathrm{erg\,cm^3\,s^{-1}}],
  \label{eq:cool}
\end{equation}%
a rate coefficient multiplying $\nelec n_{\mathrm{ion}}$, fitted over
$10^3$--$10^5$~K to the $\nelec \to 0$ limit of that ion's own
statistical-equilibrium solution --- all population in the lowest level,
excitation at the observed level energies, and the full cascade energy
radiated.  The maximum fit error over the temperatures that carry the
cooling is $2.2\%$ across the default C/N/O/Ne/S set (the coefficients of
every ion are distributed with the code, Appendix~\ref{app:atomic}; the
heavily depleted Si/Cl/Ca ions have larger full-range fit errors, confined
to the temperature tails where the cooling contribution is negligible).

By itself Eq.~(\ref{eq:cool}) carries no collisional de-excitation, so it is
exact for lines whose critical density is far above the local $\nelec$ and
increasingly too large for those below it --- the far-infrared
fine-structure transitions [O\,\textsc{iii}] 52 and 88\,$\mu$m and
[C\,\textsc{ii}] 158\,$\mu$m, whose electron critical densities are
$10^{2}$--$10^{3.5}\,\mathrm{cm^{-3}}$, so at the
$\nelec \sim 10^2\,\mathrm{cm^{-3}}$ of an \HII{} region it overestimates
the line-cooling budget by about $12\%$ (Section~\ref{sec:coolcheck}).  The
thermal loop therefore suppresses the coefficient to the local density.  A
$(\Te, \nelec)$ table, built once at setup by solving the same $n$-level
atom over a grid and stored as the ratio
$\Lambda_{\mathrm{SE}}(\Te,\nelec)/\Lambda_{\mathrm{SE}}(\Te,\nelec\!\to\!0)$,
multiplies Eq.~(\ref{eq:cool}) at each trial temperature; the fraction of the
low-density cooling the (often truncated) $n$-level model does not resolve
--- its high-excitation resonance lines, of far higher critical density ---
is left unsuppressed, and as $\nelec\to0$ the suppression is unity and the
coefficient returns to the fit exactly.  Every ion carrying a cooling fit
also carries an $n$-level model, so none is left in the zero-density limit;
the two neutrals C\,\textsc{i} and N\,\textsc{i} are the most strongly
suppressed of all, their far-infrared fine-structure lines having
$n_{\rm crit}$ of only a few $\mathrm{cm^{-3}}$.  This density-suppressed cooling
(\code{cooling\_model = local\_ne}, the default) is built from the very
solve that produces the diagnostic lines, so the two are one physics; the
pure zero-density fit is retained as the \code{low\_density} option, exact
in the diffuse gas the limit is written for, and the planetary-nebula
densities of Section~\ref{sec:summary} require the table in any case.

In neutral and partly neutral gas the ground-term fine structure of
C\,\textsc{i}, C\,\textsc{ii} and O\,\textsc{i} is excited by atomic
hydrogen as well as by electrons, and hydrogen dominates wherever
$n_{\mathrm{HI}}/\nelec \gg 1$.  These three atoms therefore carry
hydrogen-impact rates in the same $n$-level model: \citet{abrahamsson2007}
for C\,\textsc{i} and O\,\textsc{i} (their Eq.~1 with the coefficients of
their Tables~2 and 4, all three transitions of each $^3P$ term) and
\citet{barinovs2005} for C\,\textsc{ii}, each held at the end of its
temperature range rather than extrapolated, with the excitation rates
following by detailed balance.  The rate matrix of the statistical
equilibrium becomes $\nelec q^{e} + n_{\mathrm{HI}}q^{\mathrm{H}} + A$, so the
two colliders saturate the level populations together, and the emitted
lines and the thermal balance come from that one solution: the line output
includes the hydrogen impact, and the thermal loop adds to the electron
share of the $(\Te, \nelec)$ table the increment $P(e + \mathrm{H}) - P(e)$
in the radiated power.  This replaces a two-level treatment of
[C\,\textsc{ii}] 158\,$\mu$m and [O\,\textsc{i}] 63\,$\mu$m that entered
the cooling but not the emitted lines, carried no hydrogen impact for
C\,\textsc{i} at all, and used an [O\,\textsc{i}] coefficient of
undocumented origin that is 0.20--0.46 of the \citet{abrahamsson2007} rate
over 50--1000~K.

For the
output-time diagnostics, the lowest $n$ levels, radiative $A$-values, and
effective collision strengths $\Upsilon(T)$ (low-order Chebyshev fits in
Burgess--Tully scaled space; worst transition 1.7\%,
Table~\ref{tab:tier2}) are exported for the line-emitting ions of the
eleven-element registry.  A statistical-equilibrium solve at output time
(\code{nlevel\_mod}) uses these to compute the collisional emission-line
luminosities exactly --- the density-dependent emissivities the
low-density cooling fits cannot provide.  The fitted files reproduce direct
CHIANTI spline evaluations of the classic diagnostic ratios to 0.05--0.5\%,
and the Fortran solver matches an independent Python reference solver to
four significant digits.  PyNeb cross-checks give [\textsc{O\,iii}]
5007/4363 to 1.2\% and [\textsc{N\,ii}] 6584/5755 to 0.3\%; the density
diagnostics and the iron lines differ by up to 6--16\% from PyNeb's
independent collision data --- the well-known database spread, a
data-vintage bracket rather than a code error.

\subsection{Recombination lines and the nebular continuum}
\label{sec:reclines}

Hydrogen and \HeII{} recombination-line emissivities are interpolated from
the \citet{storey1995} case-A/B tables (case A when \code{par\%lyman\_lines}
is \code{'escaping'}, case B when it is \code{'on\_the\_spot'}),
and \HeI{} lines from the \citet{porter2012,porter2013} case-B
collisional-radiative emissivities.  The nebular continuum combines the
free-bound plus free-free coefficients of the MOCASSIN
\citep{ercolano2006} tables with a Milne extension above the edges and the
three two-photon continua, the \HI{} and \HeII{} ones from the share of
$\alpha_B$ that ends in $2s$ in the same level-resolved cascade
(Section~\ref{sec:recomb}), in the case the Lyman lines are given.  The
ranges those tables leave out are filled rather than left empty: the gap of
$\sim$$2\times10^{-4}$~Ryd around each edge and the stretch to the ground
edge by extending the neighboring edge pairs, He$^{+}$ longward of
1~$\mu$m, where its table ends, by the H$^{+}$ coefficient (the levels there
are hydrogenic, with the same statistical weights), and H$^{+}$ and
He$^{2+}$ beyond their tables by free-free plus Kramers capture into every
level whose edge lies below $h\nu$.  The free-free Gaunt factor is the
\citet{hummer1988} fit, with the \citet{vanhoof2014} table as a drop-in
option (the two agree to a median $5\times10^{-5}$ in $\Te$).

\subsection{The \HeI{} \texorpdfstring{$2^3$S}{2 3S} metastable level}
\label{sec:heimeta}

The metastable $2^3$S level of neutral helium is the lower level of the
10833~\AA\ triplet, whose absorption traces the extended atmospheres of
close-in exoplanets and the neutral gas of nebulae.  In \HII{} regions the
triplet is seen in emission, and near-infrared spectrographs such as
NIRSPEC, CARMENES, and X-Shooter resolve its profile in Galactic regions
like M17-B and NGC 3603 and in extragalactic star-forming regions
\citep{draine2026}.  \MoCHII{} solves its
population leaf by leaf, following the balance of \citet{oklopcic2018}:
recombination into the triplet ladder and collisional feeding from the
ground state populate $2^3$S, while radiative decay, electron-impact
transfer to the singlet levels, de-excitation to the ground state and
ionization, Penning ionization by neutral hydrogen, and photoionization by
the transported radiation field depopulate it.  The level holds a
negligible share of the helium.  Where recombination feeds it,
$n(2^3{\rm S})/n({\rm He^+})$ rises from the radiative-decay limit
$\alpha_3\nelec/A_{31} = 1.7\times10^{-7}$ at $\nelec =
100\,\mathrm{cm^{-3}}$ and $10^4$~K ($\alpha_3$ the recombination
coefficient into the triplets, $A_{31} = 1.27\times10^{-4}\,\mathrm{s^{-1}}$
the decay rate of $2^3$S) toward the high-density plateau of
$5.9\times10^{-6}$ given below, and in the dense example of
Figure~\ref{fig:hei10833} $n(2^3{\rm S})$ is at most $4.2\times10^{-6}$ of
the total helium density.  The balance is therefore linear in the
population and closes in each leaf without iteration.  From the resulting
density the line-center optical depth of the 10833~\AA\ doublet is formed.

The population is not only a diagnostic.  A recombination that ends in
$2^3$S and is then photoionized or collisionally ionized has not
recombined, and a Penning ionization ionizes a hydrogen atom; which singlet
exit a collisionally transferred atom takes decides whether its decay
photons can ionize hydrogen.  By default these processes enter the
ionization balance, the thermal balance (the excitation energy the
collisions take up and the $2^3$S photoelectrons deposit) and the exits of
the \HeI{} diffuse channel (Section~\ref{sec:diffuse}).  Since the
photoionization threshold of $2^3$S lies at 4.77~eV, the band is then
extended down to 4.78~eV, the lower end of the cross-section fit
(Section~\ref{sec:band}), unless the input turns the extension off or sets
its lower edge itself.

Several ingredients of the balance deserve note.  The rates are formed from
the effective collision strengths of \citet[][Tables~2 and 3]{bray2000}
directly, for every triplet term up to $n = 5$ from the ground state and
every term up to $n = 5$ from $2^3$S, with the exact Maxwellian constant
$h^2/[(2\pi m_{\rm e})^{3/2}k^{1/2}] = 8.629\times10^{-6}$ in cgs and the
observed level energies; the form $2.1\times10^{-8}\sqrt{13.6\,{\rm
eV}/kT}$ of \citet{oklopcic2018} is $3.3\%$ below that constant.  The
balance of \citet{oklopcic2018} carries electron-impact excitation
$1^1{\rm S}\to2^3{\rm S}$ but not its inverse, without which the two
levels cannot approach their Boltzmann ratio; \MoCHII{} adds the
de-excitation $2^3{\rm S}\to1^1{\rm S}$ by detailed balance.  It is
$5.4\%$ of the electron-impact transfer to the singlets at $10^4$~K and
$56\%$ at 3000~K.  Recombination into the triplets is the triplet share of
the same \HeI{} cascade the diffuse channel uses times the case-B
coefficient of the ionization balance, rather than a separate fit that
would normalize the same recombinations a second time.  Electron-impact
ionization of $2^3$S, absent from that balance, takes
$8.38\times10^{-10}\,T^{1/2}\exp(-55338/T)\,\cmc$
\citep[][Eq.~12]{black1981}; it is $0.97\%$ of the transfer to the
singlets at $10^4$~K and $18\%$ at $2\times10^4$~K.  The Penning and associative
ionization by neutral hydrogen takes the rate coefficient of
\citet[][Fig.~5]{garciamunoz2025}, a Maxwellian average and therefore
continuous in $T$; the two-branch power law of \citet{taylor2025} that it
replaces jumps by a factor of $1.57$ at 4000~K.  The photoionization rate
of $2^3$S is accumulated during transport as a path estimator at each
packet's own energy, as the other photoionization rates are, because its
cross section has structure inside a diagnostic bin (a Cooper minimum near
35~eV).  At $10^4$~K the high-density plateau of
$n(2^3{\rm S})/n(\mathrm{He^{+}})$, where recombination into the triplets is
balanced by electron-impact removal, is $5.867\times10^{-6}$, and the
electron density at which it meets the radiative-decay limit is
$3502\,\mathrm{cm^{-3}}$.

Figure~\ref{fig:hei10833} shows both sides of the transition for a dense
compact example ($\nH = 10^4\,\mathrm{cm^{-3}}$, $R = 0.15$~pc, $\Te$ fixed
at $10^4$~K).  On the absorption side, the $2^3$S column density reaches
$N = 2.87\times10^{15}\,\mathrm{cm^{-2}}$ and the line-center optical depth of
the blended $J' = 1, 2$ doublet reaches $3.47\times10^{3}$ ($2.98\times10^{3}$
on the central sightline), falling through $\tau_0 = 1$ only at a projected
radius of $0.146$~pc, so all but the outermost rim of the projected disk is optically
thick at 10833~\AA.  Such a column makes the line resonantly scatter, as
\citet{draine2026} shows for \HII{} regions: repeated scattering broadens the
emergent profile into a multipeaked feature whose width can exceed
$100\,\mathrm{km\,s^{-1}}$ and whose centroid is blueshifted by up to
$\sim$$14\,\mathrm{km\,s^{-1}}$ relative to other lines, while the extended
path length lets dust absorb a larger share of the photons and so weakens the
line.  The He\,\textsc{i} 10833/Pa$\gamma$ ratio therefore stops being a clean
He$^{+}$/H$^{+}$ measure once the metastable column is large.  The
optically-thin diagnostic likewise underestimates the true absorption over
most of the region, which is what makes the $2^3$S output of
Figure~\ref{fig:hei10833} the input to a scattering calculation rather than an
observable in itself.  That calculation is within reach: the converged $2^3$S
density is handed to \textsc{LaRT}
\citep{seon2020,seon2022,seon2024}, a Monte Carlo resonant-line transfer code
that already carries a 10833~\AA\ module and has been applied to the line in
an exoplanet atmosphere \citep{yan2022}.  The coupling is in place --- a
converter writes the \MoCHII{} metastable grid in \textsc{LaRT}'s input
format, and the two codes agree on the line-center optical depth of the
example above to $0.5\%$, measured on the run made before the atomic-data
update of Section~\ref{sec:atomic} --- and the resonant-scattering profiles will be
presented in a forthcoming paper.  On the
emission side, the same
transition radiates the 10833~\AA\ line, and the code writes its emissivity
from the \citet{porter2012} recombination table --- a collisional-radiative
total, so it is used as tabulated and nothing is added to it.  Integrated, the
line carries $1.33\,L(\mathrm{H}\beta)$ here, and its emissivity falls
inside the ionized sphere: the He$^{+}$ fraction falls through one half at
$r = 0.138$~pc, where the $>$24.6-eV photons run out, while hydrogen stays
ionized to $0.145$~pc.

\begin{figure*}[t]
\centering
\includegraphics[width=0.95\textwidth]{figures/hei10833.pdf}
% source: python/make_hei10833.py, from
%   results/continuous_energy/m11_hecx_solver/hei_high_density_converged_{heimeta,emis}.h5
%   (examples/hei_metastable/hei_high_density.in, 2026-09-25 defaults, hydrogenic
%   alpha_B and the stopping test of 2026-09-26, the He+ + H charge exchange and
%   hydrogen-balance solve of 2026-09-27)
\caption{The \HeI{} $2^3$S metastable level for a dense compact
example ($\nH = 10^4\,\mathrm{cm^{-3}}$, $R = 0.15$~pc, $\Te = 10^4$~K), as
radial profiles of the two sides of the same $2^3{\rm S}$--$2^3{\rm P}$
transition.  Left, absorption: the 10833~\AA\ line-center optical depth
$\tau_0$ of the blended $J' = 1, 2$ doublet against projected radius (left
axis) with the $2^3$S column density on the right axis.  Both are
line-of-sight integrals, and at the fixed $\Te$ of this run the Doppler width
is constant, so the two are strictly proportional and the right axis is the
conversion of the left; the curves overlie.  The region is optically thick
inside $b = 0.146$~pc, where the optically-thin diagnostic breaks down.
Right, emission: the 10833~\AA\ emissivity against radius, from the
\citet{porter2012} He\,\textsc{i} block of the emissivity file.}
\label{fig:hei10833}
\end{figure*}

%==========================================================================
\section{Validation}
\label{sec:validation}

\MoCHII{} is validated against analytic solutions, an exact 1D reference
solver, and the Lexington benchmark suite, in staged tests.  Each
AMR result is required to reproduce the equivalent uniform-grid run.

\subsection{Rate integrals and the Str\"omgren sphere}
\label{sec:stromgren}

The zero-variance edge walk (Eq.~\ref{eq:edgewalk}) is verified on a
uniform sphere ($\nH = 1$, $R = 1$~pc, $\xHI = 0.1$) with a central
$10^5$~K source.  Against the cell-volume-averaged analytic attenuation,
$\Gamma_{\mathrm{HI}}$ has a mean bias of $+0.08\%$ and a median deviation
of $0.40\%$ (pure Monte Carlo noise), and $\Gamma_{\mathrm{HeI}}$ a bias of
$+0.003\%$ (Figure~\ref{fig:g0}).

\begin{figure}[t]
\centering
\includegraphics[width=\columnwidth]{figures/g0_gamma_check.pdf}
% source: python/make_g0_gamma.py, from tests/g0_gamma/g0_gamma_rates.h5
\caption{Rate-integral test: $\Gamma(r)$ from the Monte Carlo tally
(points) against the analytic point-source attenuation (lines); the lower
panel shows the ratio.  Shaded regions are excluded (cell-size effects at
the center, the voxelized sphere edge outside).}
\label{fig:g0}
\end{figure}

The Str\"omgren test uses a uniform $\nH = 100$ sphere, $\QH =
10^{49}\,\mathrm{s^{-1}}$ ($T_\star = 4\times10^4$~K), case B, and
$\Te = 10^4$~K.  For the reference we solve the spherically symmetric
ionization-equilibrium equations directly, with a deterministic radial
marching over $2\times10^4$ shells that carries the same band bins,
photoionization cross sections, and recombination and collisional rates as
the code: each bin's photon flux is attenuated with the running optical
depth evaluated at mid-shell, and the H/He ionization equilibrium is
solved shell by shell.
The sweeps must update opacities Gauss--Seidel style (a Jacobi sweep,
using the previous sweep's opacities, advances the front only one
absorption length in each pass); a few sweeps then converge.  ``Exact'' means
free of Monte Carlo noise and grid-converged at this resolution --- in
the sharp-edge limit it reproduces the analytic Str\"omgren radius
$R_S = [3\QH/(4\pi \nH^2 \alpha_B)]^{1/3} = 3.157$~pc, from which the
full solution differs physically (helium consumes photons; the front has
a finite width).  The 1D reference
gives an effective ionized radius $R_{\mathrm{eff}} = (3V_{\mathrm{ion}}/
4\pi)^{1/3} = 3.0355$~pc; a uniform level-7 octree gives 3.0563~pc
($+0.69\%$) and a level-6 one 3.0663~pc ($+1.01\%$), the difference being the
finite front width on a 0.125~pc cell against a front at 3.04~pc; a
front-refined octree with $3.5\times$ fewer cells
reproduces the level-6 radius to $5\times10^{-4}$~pc (3.0668~pc, $+1.03\%$), while a center-refined control (coarsest
cells at the front) deviates by $+2.80\%$ --- the ionization front is the
one place needing resolution (Figure~\ref{fig:g1}).  With the thermal
balance on, $\Te(r)$ matches the 1D thermal reference to a median
$0.12\%$, a mean bias of $-0.004\%$ and a largest deviation of $3.4\%$
(Figure~\ref{fig:g2a}).  The reference carries every term the run does:
besides the recombination, free--free and cross-section routines of the
code, taken through a shared module that a gate holds to the Fortran to
machine precision, it carries the hydrogen share of the \HeI{} and \HeII{}
on-the-spot photons (Section~\ref{sec:diffuse}), the $2^3$S coupling
(Section~\ref{sec:heimeta}) and the band below the hydrogen edge down to
4.78~eV; only the helium--hydrogen charge exchange of Section~\ref{sec:ci}
is switched off in the run, since the reference does not carry it.  Its
stellar band is integrated on 256 bins, converged to $5\times10^{-5}$ in
$\Te$: the 32 bins used earlier put $\Te$ low by $0.42\%$ on average, which
was most of the mean bias of $+0.61\%$ reported before, and aligning the
recombination data with the run had moved the median from $0.97\%$ to
$0.77\%$.  The run stops on the drift test of Section~\ref{sec:iterate},
restarted from the previous converged state; the single-cell maxima never
settle (they reached the 40-iteration cap in the earlier runs), while the
volume-integrated change of $\Te$ sits at its Monte Carlo floor of
$0.22\%$.  Inside $r = 0.24$~pc both the run and the reference reach the
$\Te$ ceiling of $5\times10^4$~K: there the \HeII{} recombination cascade
absorbed on the spot heats gas that is optically thin to it (the optical
depth from the center to 0.24~pc is $2\times10^{-6}$ at 40.8~eV), which the
on-the-spot approximation does not describe; the comparison region starts
outside it, at 0.3~pc.  Transporting the recombination continuum instead
holds that core at about $3.1\times10^4$~K.  What is left in the comparison
region is a property of the solution rather than of the data.  Without metals the
metal-free nebula runs at 16--27~kK, the known overshoot that motivates the
trace-metal cooling.

\begin{figure*}[t]
\centering
\includegraphics[width=0.9\textwidth]{figures/g1_stromgren_check.pdf}
% source: python/make_g1_stromgren.py, from
%   results/continuous_energy/m10_drift_test/g1/ (inputs tests/g1_stromgren/)
\caption{Str\"omgren test: H and He ionization structure against the 1D
exact reference (black); uniform level-6 (blue), center-refined (red),
and front-refined (green, overlying blue).}
\label{fig:g1}
\end{figure*}

\begin{figure*}[t]
\centering
\includegraphics[width=0.9\textwidth]{figures/g2a_thermal_check.pdf}
% source: python/make_g2a_thermal.py, from
%   results/continuous_energy/m11_hecx_solver/g2a_thermal_rates.h5
\caption{Thermal test: $\Te(r)$ and $\xHI(r)$ with the thermal balance,
against the 1D exact reference.}
\label{fig:g2a}
\end{figure*}

\subsection{The Lexington benchmarks}
\label{sec:hii}

The HII20 and HII40 cases of the Lexington/Meudon comparison
\citep{ferland1995,pequignot2001} use 20 and 40~kK blackbody sources
($\QH = 1.0$ and $4.3\times10^{49}\,\mathrm{s^{-1}}$), $\nH = 100$ in a
hollow shell, and the standard Lexington abundances (He 0.1, C
$2.2\times10^{-4}$, N $4\times10^{-5}$, O $3.3\times10^{-4}$, Ne
$5\times10^{-5}$, S $9\times10^{-6}$).  They are the natural reference for a
new code because eight independent codes published results for the same two
models, so the benchmark carries its own measure of how large a
disagreement is significant.  Tables~\ref{tab:wood40} and~\ref{tab:wood20}
place \MoCHII{} next to those eight, in the compilation of
\citet{ercolano2003}, together with the Monte Carlo code of
\citet{wood2004} that was added to the same tables afterwards.

The \MoCHII{} runs compared here are configured to match the published
calculations wherever the setup is a choice rather than a physical
difference.  The ionizing luminosity is set from the benchmark blackbody
luminosity so that $\QH$ reproduces the specified $4.265\times10^{49}$
(HII40) and $1.004\times10^{49}\,\mathrm{s^{-1}}$ (HII20).  Stellar photon
energies are sampled continuously from the Planck spectrum with a Sobol
sequence, and every photoionization cross section is evaluated at the
sampled energy (Section~\ref{sec:band}); the 32 band bins are only a
diagnostic tally, not kept in these runs.  The band extends below the hydrogen edge to 4.78~eV,
just above the 4.77~eV photoionization threshold of \HeI{} $2^3$S (the lower
end of its cross-section fit), whose population the ionization and
thermal balances carry (Section~\ref{sec:heimeta}); those photons also
ionize the neutral stages of carbon and sulfur, as the full stellar spectrum
does in \textsc{Cloudy}.  The diffuse field is transported explicitly with
case-A rates and the \HeI{} and \HeII{} channels
(Section~\ref{sec:diffuse}), the treatment the published codes use, so that
the comparison is not confounded by the transfer treatment.
The recombination-line emissivities
stay on the case-B tables, because the Lyman lines are optically thick in
these models whatever the continuum treatment is; the two cases are
separate switches (\code{par\%recombination\_continuum} for the ionization,
\code{par\%lyman\_lines} for the line tables).  The grid is a uniform level-6
octree ($64^3$ leaves) run with $2^{25}$ stellar packets in the ionizing band
and $2^{23}$ below it (Sobol seed 20260728) to the drift test of
Section~\ref{sec:iterate}.  The runs reached it after 64 (HII40) and 55
(HII20) iterations, and were then restarted from those states with the
charge-exchange heating, the $2^3$S coupling, the band below the hydrogen
edge, the \HeII{} diffuse channel and the collider-resolved
two-photon quenching switched on, reaching it again after 8 and 18 further
iterations.  They were restarted twice more: with the hydrogenic $\alpha_B$
of Section~\ref{sec:recomb}, continued to the stopping test in its present
form (44 and 41 further iterations), and with the helium--hydrogen charge
exchange of Section~\ref{sec:ci} and the hydrogen balance solved for
$\ln(n_{\mathrm{HII}}/n_{\mathrm{HI}})$ (10 further iterations each).  At HII40 the gas extends to
5.2~pc ($1.60\times10^{19}$~cm) rather than to the $1.46\times10^{19}$~cm
outer radius of the \textsc{Mocassin} benchmark input: run to convergence,
the hydrogen front sits at that radius (as it does in \textsc{Cloudy}, whose
$x(\mathrm{H^{0}}) = 0.5$ falls at $1.461\times10^{19}$~cm), and a grid that
ends there cuts the partly neutral edge.  At HII20 the benchmark's
$0.95\times10^{19}$~cm lies $0.2$~pc beyond the front and is kept.

We also ran \textsc{Cloudy} c25.00 \citep{cloudy2025} on both benchmarks and
carry it as the C25 column of the same two tables, providing a uniform
one-dimensional reference alongside the published multi-code spread.  Its
physics input is the
unmodified regression deck that \textsc{Cloudy} distributes for each case,
\code{tsuite/auto/hii\_paris.in} for HII40 and \code{hii\_coolstar.in} for
HII20 --- the same source, density, inner radius, and Lexington abundances as
above --- with only output commands appended; line luminosities are taken
from a \code{save line list absolute} file and divided by the \textsc{Cloudy}
H$\beta$ luminosity, and the averaged temperature is \textsc{Cloudy}'s own
$n_{\mathrm{e}}n_{\mathrm{p}}$-weighted volume average, the same weighting
used on both other sides.  The one geometric difference is that each deck
runs to its own stopping criterion rather than to a prescribed outer radius
--- an electron fraction of $10^{-2}$ at HII40, the 1000~K temperature floor
at HII20 --- reaching $1.467\times10^{19}$ and $8.97\times10^{18}$~cm, at
the top of the published $R_{\mathrm{out}}$ ranges (1.45--$1.47\times10^{19}$
and 8.7--$9.01\times10^{18}$~cm).  The \MoCHII{} gas extends beyond both, so
the two carry the same edge to within where \textsc{Cloudy} stops; the
low-ionization rows fed by that edge --- [O\,\textsc{i}], [S\,\textsc{ii}],
[N\,\textsc{ii}] --- are the ones where a remaining difference in the outer
boundary would show.

The temperature compared throughout this section is
$\langle T[N_{\mathrm{p}}N_{\mathrm{e}}]\rangle$, the electron temperature
averaged with the weight $n_{\mathrm{HII}}\nelec V$ --- the same definition
on both sides, evaluated leaf by leaf in \MoCHII{} and cell by cell in the
published models.  The inner-edge temperature $T_{\mathrm{inner}}$ that the
published tables also list is quoted separately where it is used.  Ionic
fractions are $\nelec V$-weighted means, the convention of MOCASSIN's own
output.  As a further point of comparison we also ran MOCASSIN 2.02 itself,
converged, on the HII40 model, and quote it below alongside the published
columns.

\emph{Ionization structure.}  The helium and metal ionization is sensitive
to the photon distribution around the \HeI{} edge, so the calculations
sample the Planck probability density continuously and use the cross
section at each sampled energy.  The resulting
$\langle\mathrm{He^{+}}\rangle/\langle\mathrm{H^{+}}\rangle$ is 0.752 at
HII40, inside the published 0.715--0.829 (median 0.767) and $1.5\%$ above
C25's 0.741, and 0.047 at HII20, next to the published median of 0.048
inside a spread of 0.034--0.077.

\emph{Temperature.}  The weighted temperature is 8180~K at HII40 and
6940~K at HII20.  Both match the \textsc{Cloudy} c25.00 solution of the
same two models --- 8210~K and 6910~K --- to $0.4\%$ ($-0.37\%$ and
$+0.44\%$).  At HII40 \MoCHII{} lies at the top of the 2000-era published
range (7720--8199~K, median 8026) and \textsc{Cloudy} just above it; at
HII20 both lie above it (6402--6749~K, median 6662).  The inner-edge
temperature is inside the published spread at both, 7525~K at HII40 against
a published 7370--8318~K and 6947~K at HII20 against a published
6562--7224~K.  The converged
MOCASSIN 2.02 run of the HII40 model gives 8176~K, inside the published
spread.

The temperature is set by the cooling treatment, and the two limits of that
treatment straddle the published medians.  The
density-suppressed cooling used above --- the level populations at the local
$\nelec$, Section~\ref{sec:cooldata} --- gives the 8180 and 6940~K that
track \textsc{Cloudy}; the zero-density limit of the same fits, retained as
an option, overestimates the far-infrared line cooling at
$\nelec \sim 10^{2}\,\mathrm{cm^{-3}}$ and so places \MoCHII{} instead at
7877 and 6282~K --- inside the published range at HII40 and 120~K below
its lower edge at HII20 --- which is $4.1\%$ and $9.1\%$ below
\textsc{Cloudy} c25.00.  The offset between the two limits agrees in
sign and size with the far-infrared overestimate, which the next subsection
measures directly.

% Published columns: Tables 4 and 5 of Ercolano et al. (2003), transcribed
% and checked value by value against the paper.
% MoCHII column: continuous-energy benchmark runs (Sobol seed 20260728,
% explicit diffuse field with the Milne ground continua, Q(H) matched to the
% specification, 2^25 ionizing + 2^23 FUV packets, conv_crit='vol' with the drift stopping test,
% HII40 gas extended to 5.2 pc, two-way charge exchange as the
% default, the corrected oxygen charge-exchange data, the 2026-09-24
% atomic-data update, and the 2026-09-25 defaults: charge-exchange heating,
% the He I 2^3S coupling with the FUV band from 4.78 eV, collider-resolved
% two-photon quenching, the He II diffuse channel; restarted from the
% m5_converged states, then with the hydrogenic alpha_B of 2026-09-26 (m9)
% and continued to the stopping test revised 2026-09-26 (m10), then with the
% He+ + H charge exchange and hydrogen-balance solve of 2026-09-27 (m11));
% line entries are
% the L/L(Hbeta) column of
% results/continuous_energy/m11_hecx_solver/hii{40,20}_converged_lines.txt
% (doublets summed), temperatures are from the matching rates files.
% C25 column: tests/cloudy_c25/, Cloudy c25.00 on the distributed decks.

\subsection{Cooling-budget verification at a fixed state}
\label{sec:coolcheck}

An emission-line comparison cannot by itself separate an error in the
atomic data from an error in the temperature, because each code's
temperature is the root of its own energy balance: too little cooling in
one channel raises $\Te$, which raises every other channel, and the
converged line ratios absorb part of the discrepancy.  The cooling data can
be compared free of that feedback only at a \emph{fixed} state.
\textsc{Cloudy} c25.00 writes such a decomposition --- its
\code{save cooling each} output lists the cooling rate of every ion and
every continuum channel zone by zone --- so the test is to take its
converged $(\Te, \nelec$, ionic fractions$)$ zone by zone, evaluate the
\MoCHII{} cooling terms on exactly that state, and integrate both over the
model volume.  What this measures is the absolute scale of two independent
atomic-data compilations, channel by channel, with nothing else in either
code entering.  It is the comparison that isolates the cooling data from
the temperature, which an equilibrium comparison cannot do, since the
latter shows only the temperature at which each code's own heating and
cooling cross.

On the HII40 model the two agree within the accuracy of the underlying
collision data.  The comparison is reproducible with
\code{tests/cloudy\_c25/run\_fixed\_state.sh}, which also re-evaluates the
superseded treatments and recovers the values recorded for them.  With the
$n$-level emissivities evaluated at the local $\nelec$, and with the
atomic-hydrogen impact of Section~\ref{sec:cooldata} in the same
statistical equilibrium, the ratio \MoCHII/\textsc{Cloudy} is 1.02 for [O\,\textsc{iii}],
0.98 for [O\,\textsc{ii}], 1.00 for [S\,\textsc{iii}], 0.99 for
[N\,\textsc{ii}], 1.00 for [S\,\textsc{ii}], 1.00 for [Ne\,\textsc{ii}],
1.02 for [N\,\textsc{iii}], 1.02 for [S\,\textsc{iv}], 0.99 for
C\,\textsc{iii}], 0.94 for [C\,\textsc{ii}] and 0.995 for the \HI{}
collisional lines; summed over all line channels the two agree to $1.1\%$.
Two channels fall outside $\pm$10\%, [Ne\,\textsc{iii}] (1.30) and
[O\,\textsc{i}] (0.87, of which the hydrogen impact contributes the rise from
0.84 with electrons alone), and they carry $1.9\%$ and $0.1\%$ of the
budget.
The free-free continuum, an independent
check of the Gaunt factor and of the net ion charges entering it, agrees to
$0.2\%$.  Recombination cooling differs by bookkeeping rather than by data:
\textsc{Cloudy} reports the free-bound loss net of the escaping
recombination continuum, $4.02\times10^{37}\,\mathrm{erg\,s^{-1}}$, which
falls between the \MoCHII{} case-A and case-B sums ($6.43$ and
$3.57\times10^{37}$, from Eq.~\ref{eq:betarec}), the difference being the
ground-state term that the explicit diffuse field re-emits instead of
losing.  The heating side of the
same \textsc{Cloudy} model balances its cooling to $1.3\times10^{-4}$ and
deposits 4.57~eV for each ionizing photon, so the comparison state is
itself an equilibrium one.

The same $n$-level emissivities at the local $\nelec$ are what the thermal
loop uses by default (Section~\ref{sec:cooldata}), so this fixed-state
agreement is the cooling the converged benchmark carries.  Replacing them
by the zero-density limit of the same fits (Eq.~\ref{eq:cool}, the
\code{cooling\_model = low\_density} option) raises the summed line cooling
to $1.116$ times \textsc{Cloudy}'s.  The excess is not spread over the
ions: it is 5.1 percentage points of the line budget in [O\,\textsc{iii}]
and 3.9 in [C\,\textsc{ii}], and in both it sits entirely in the
far-infrared fine-structure lines, whose critical densities are at or below
the local $\nelec = 110\,\mathrm{cm^{-3}}$.  Since the line channels carry
$79\%$ of the cooling budget, that is a $9\%$ excess in the total, and it
accounts for the temperature difference quantitatively: scaling $\Te$ on the
fixed state until the \MoCHII{} zero-density terms radiate the same total
power --- \textsc{Cloudy} balances its own heating at the unscaled
temperature, and the photoheating of the two codes agrees --- gives
7792~K against \textsc{Cloudy}'s 8210~K.  Here the \MoCHII{} total is the
zero-density line cooling plus the case-B recombination cooling of
Eq.~(\ref{eq:betarec}) plus the remaining continuum terms, all evaluated at
$f\,T$ zone by zone, and 7792~K is $f$ times \textsc{Cloudy}'s
$T(N_{\mathrm{e}}N_{\mathrm{p}}) = 8210$~K; it is close to the 7877~K
the converged zero-density run reaches.  Removing the overestimate, the
density-suppressed default reaches 8180~K, matching \textsc{Cloudy} to
$0.4\%$ (Section~\ref{sec:hii}).  The far-infrared overestimate, not the atomic
data, is thus what separates the two limits, and the $(\Te,\nelec)$ table of
Section~\ref{sec:cooldata} removes it.

\emph{Sensitivity to the other ingredients.}  The remaining modeling
choices are smaller levers than the cooling function, and each is measured
by an in-code experiment.  The transport treatment is the largest of them:
replacing the case-A explicit diffuse field by the case-B on-the-spot
approximation, with every other setting fixed, gives 8129~K at HII40 and
6915~K at HII20, so the explicit field used here raises $\Te$ by
$0.6\%$ and $0.4\%$ and moves the two benchmarks the same way.  That
distance depends on what the on-the-spot approximation does with the
\HeI{} recombination photons.  Given back to helium, as it was before the
treatment of Section~\ref{sec:diffuse} gave hydrogen its share, it left
HII40 at 7962~K, which the explicit field exceeded by $3.1\%$, with the
hydrogen front $0.18$~pc inside the explicit field's; with hydrogen's share the front is $0.015$~pc inside.  Substituting
MOCASSIN-style recombination
(P\'equignot 1991 radiative plus Nussbaumer--Storey 1983 dielectronic) for
the Badnell 2023 data, whose $\alpha(\mathrm{X\,\textsc{iii}} \to
\mathrm{X\,\textsc{ii}})$ are systematically lower (ratios 0.72--0.84),
moved $\Te$ by $0.15\%$ and the carbon, nitrogen and oxygen ionization
fractions by about 0.02 in runs stopped by the earlier convergence rule of
Section~\ref{sec:iterate}, a difference at the resolution of that rule, so
the recombination vintage is at most a minor lever on both.  The dominant
[O\,\textsc{iii}] collision strengths are stable across database
generations: the $^3\mathrm{P}$--$^1\mathrm{D}$ effective collision
strength is 2.283 in the CHIANTI~v7 compilation behind the published
MOCASSIN column and 2.262 in v11 at $10^4$~K, a $0.9\%$ difference worth
less than 30~K.  The photoionization heating is standard: weighted by the
photoionization rate over the HII40 nebula, the mean photoelectron energies
are $\langle h\nu - 13.6\rangle_{\mathrm{H}} = 2.9$~eV and
$\langle h\nu - 24.6\rangle_{\mathrm{He}} = 3.2$~eV, against 3.0 and
3.5~eV for the unattenuated 40~kK blackbody.  Grid resolution
is treated separately below.

\emph{Line ratios.}  The line intensities follow from that temperature and
ionization structure rather than from the emission calculation, which is
exact to the accuracy of its input data (Section~\ref{sec:cooldata}: the
$n$-level solver reproduces an independent reference solver to four digits,
and the fitted collision data reproduce direct CHIANTI evaluations to
0.05--0.5\%).  At HII40 the infrared fine-structure lines, whose excitation
energies are small enough that they are nearly temperature-independent, sit
with the published set ([Ne\,\textsc{iii}] 15.5\,$\mu$m 0.282 in a
published 0.263--0.429; [Ne\,\textsc{ii}] 12.8\,$\mu$m 0.206 in
0.177--0.217; [O\,\textsc{iii}] 52+88\,$\mu$m 2.45 just above 2.10--2.42),
and so do the temperature-sensitive optical lines of the same ion:
[O\,\textsc{iii}] 5007+4959 is 2.36 against a published 1.64--3.27, and the
auroral [O\,\textsc{iii}] 4363, the most temperature-sensitive entry in the
tables, is 0.0044 against a published 0.0022--0.0070.  The lines that fall
outside the published ranges are examined below.

\emph{H$\beta$ and the photon budget.}  A recombination line is a photon
counter, and $L(\mathrm{H}\beta)$ therefore tests whether the model is
radiation-bounded: every ionizing photon absorbed by hydrogen returns as a
recombination, and a fixed fraction of those recombinations produces
H$\beta$.  Both models here are bounded.  At HII40 the shell-averaged
$x_{\mathrm{HII}}$ falls through 0.5 at $1.464\times10^{19}$~cm, inside gas
that extends to $1.60\times10^{19}$~cm; at HII20 the outer shells are fully
neutral.  The budget closes accordingly: summing the case-B recombinations
over the converged solutions gives $0.98\,\QH$ and $1.00\,\QH$, the rest of
the HII40 photons being taken up by helium, most of whose energy returns to
hydrogen through the \HeI{} diffuse channels.  The predicted
$L(\mathrm{H}\beta) = 2.02\times10^{37}$ and
$4.85\times10^{36}\,\mathrm{erg\,s^{-1}}$ are $1.3\%$ and $0.9\%$ below the
published medians and inside the published ranges (2.01--2.10 and
4.83--5.04); against C25 they are $1.7\%$ and $2.9\%$ low.  Before the
stopping rule of Section~\ref{sec:iterate} was replaced they were
$4$--$6\%$ below the medians, the early stop leaving the H\,\textsc{ii}
region a few percent short in volume, and the floor test of that section
added the last $0.7\%$ at HII40.

\emph{The complete published line list.}  The full benchmark line lists,
with all published codes, a \textsc{Cloudy}~c25.00 column, and the
\MoCHII{} column, are given in Appendix~\ref{app:lines}
(Tables~\ref{tab:wood40} and~\ref{tab:wood20}); they repeat the comparison
over the full benchmark line list in
the form used by \citet{wood2004}, whose Tables 1 and 2 add their own Monte
Carlo code (WME) to the eight other codes of the comparison; the
\MoCHII{} and C25 columns provide two additional reference solutions.  The
low-ionization edge line [O\,\textsc{i}]
6300+6363 is 0.0281 at HII40, above the published 0.0059--0.012 range; it is
particularly sensitive to the thin partly neutral edge, and it is the one
entry the corrected oxygen charge exchange of Section~\ref{sec:ci} moves
appreciably ($+39\%$, by raising the neutral-oxygen fraction at the edge,
against $+0.03\%$ in the weighted temperature), as does the band below the
hydrogen edge, which heats that edge ($+14\%$).  The complete
comparison is given directly in the tables.

The C25 column provides a uniform reference, and with the two temperatures
matched the line
ratios track C25 across the board.  At HII40, where the weighted temperature
is within $0.4\%$ of C25's, both the temperature-insensitive far-infrared lines and the
temperature-sensitive optical and auroral lines agree with C25 to within
$6\%$: among the infrared lines [O\,\textsc{iii}] 52+88\,$\mu$m is
$1.02\times$ C25, [N\,\textsc{iii}] 57.3\,$\mu$m $1.01\times$ and
[S\,\textsc{iii}] 18.7\,$\mu$m $1.00\times$; among the optical ones
[O\,\textsc{iii}] 5007+4959 is $0.96\times$, its auroral 4363 $0.94\times$
and [O\,\textsc{ii}] 3726+3729 $0.98\times$ --- the same optical and auroral
lines that the zero-density cooling, running $3.7\%$ cooler, suppresses
to $0.68$--$0.82\times$.  That the temperature-sensitive lines fall into
line with C25 exactly when the temperature does is the line-level
counterpart of the cooling match.  At HII20, where \MoCHII{} is $0.4\%$
hotter than C25, [O\,\textsc{ii}] 3726+3729 is $1.02\times$ C25.  The
residual outliers at HII40 are not temperature: [Ne\,\textsc{iii}]
15.5\,$\mu$m ($1.30\times$) and 3869+3968 ($1.21\times$) and
[O\,\textsc{i}] 6300+6363 ($2.51\times$, on a line the thin neutral edge and
the grid resolution govern), each an ionization or collision-strength
difference in a single ion.  \HeI{} 5876, which the early-stopped runs put
at 0.107, below the published 0.112--0.119 at HII40, is 0.112, at its lower edge.

The C25 and published columns illustrate the benchmark's code-to-code
spread.  [S\,\textsc{ii}] 6716+6731 at
HII40 is the clearest case: the GF entry is 0.137, C25 gives 0.304 and
\MoCHII{} 0.360; the two latter values differ by $18\%$.  Sulfur at HII20 is
the same story in the other direction: [S\,\textsc{iii}] 18.7\,$\mu$m
(0.212 against a published 0.34--0.46) and [S\,\textsc{iii}] 9532+9069
(0.299 against 0.465--0.604) fall \emph{below} every published code, and so
does C25 (0.209 and 0.293), while [S\,\textsc{ii}] 6716+6731 sits inside the
published range for both (0.806 and 0.744).  Both C25 and \MoCHII{} have the
same S\,\textsc{ii}/S\,\textsc{iii} ordering, which points at the sulfur
data rather than at either code.  The band below the hydrogen edge closes
one further difference: \textsc{Cloudy} carries the full stellar spectrum,
so photons below 13.6~eV ionize C$^{0}$ (11.26~eV) and S$^{0}$ (10.36~eV),
and a band starting at the hydrogen edge leaves both neutral wherever
hydrogen is.  Compared in 0.05~pc shells (volume-weighted over the gas
leaves) with \textsc{Cloudy}'s zones at the same radius, the \MoCHII{}
C$^{0}$ fraction without those photons was 1.7--4.5 times
\textsc{Cloudy}'s between 1.0~pc and the HII40 \HeI{} front (4.27~pc) and
3.1--3.7 times between 1.0~pc and HII20's (1.25~pc), and between the
\HeI{} and H fronts (to 4.70 and 2.88~pc) it rose from 4.4 to 190 times
(HII40) and from 3.2 to 540 times (HII20); it averaged $5.7\%$ over the
4.4--4.75~pc shell that holds the HII40 front.  With them, in the run of
the tables, it is 0.99--1.33
and 1.06--1.19 times \textsc{Cloudy}'s over the same two HII40 intervals and
1.00--1.16 and 0.97--1.05 times over the HII20 ones ($4.5\times10^{-4}$ in
that shell, against \textsc{Cloudy}'s $4.1\times10^{-4}$), and
[C\,\textsc{i}] 9824+9850 falls by a factor of 65 at HII40 and 90 at HII20.
The same photons heat the partly neutral edge, by up to 900~K (HII40) and
750~K (HII20) in the 0.05~pc shell just outside the front, which raises [O\,\textsc{i}]
6300+6363 by $14\%$ and $22\%$ and [N\,\textsc{i}] 5198+5200 by $20\%$
and $25\%$, while [S\,\textsc{ii}] 6716+6731 rises by $4\%$ and $2\%$ and
[S\,\textsc{iii}] does not move --- so the band is not the cause of the
[S\,\textsc{iii}] difference either.  Charge-exchange heating, the $2^3$S
coupling in the ionizing band and the two-photon quenching, restarted from
the same states without the band, move no line by more than $1.6\%$, and
the \HeII{} diffuse channel, which that restart did not include, has almost
nothing to carry here: the volume-averaged He$^{2+}$ fraction is
$9\times10^{-5}$ at HII40 and $6\times10^{-9}$ at HII20.  These changes are
therefore the band's.

A second, structural difference is worth naming here, because it affects
sulfur specifically.  The published columns are not all built on the same
set of ionization stages: the code of \citet{wood2004}, and
\textsc{CMacIonize} after it, omit neutral sulfur and start the sulfur
chain at S$^{+}$, so that in gas where the ionizing field has been absorbed
all sulfur is S$^{+}$ by construction; oxygen and neon in the same codes do
carry their neutral stages.  \textsc{Mocassin} instead solves the neutral
stage of every element as an unknown, as \MoCHII{} and \textsc{Cloudy} do.
The convention matters only outside the ionization front, where it fixes
the S$^{+}$ reservoir that feeds [S\,\textsc{ii}] rather than letting it
recombine to S$^{0}$; it is one of the ingredients of the scatter among the
published sulfur entries, and it is worth keeping in view when a single
[S\,\textsc{ii}] ratio is used to judge a code.

\emph{Radial structure.}  Figures~\ref{fig:wood40} and \ref{fig:wood20}
give the radial ionization and temperature structure in the panel layout of
\citet{wood2004}, their Figures 1 and 2, so that the two can be read panel
by panel.  The morphology is the published one throughout.  At HII40 the
hydrogen and helium neutral fractions rise monotonically from
$4\times10^{-5}$ and $2\times10^{-4}$ at the inner edge, the singly and
doubly ionized fractions of carbon, nitrogen and oxygen cross within
0.3~pc of one another, between 3.9 and 4.2~pc, C$^{+3}$ peaks at $0.05$ and is
confined to the inner region, and the temperature falls slightly out to
about 2.3~pc before climbing steeply toward the front.  At HII20 helium is
essentially all neutral beyond
1.3~pc while hydrogen stays ionized out to 2.9~pc, so the doubly ionized
metals occupy only a thin inner shell.  Both models are radiation-bounded,
so both outer radii are predictions.  The ionized-volume-equivalent radius
$(3V_{\mathrm{ion}}/4\pi)^{1/3}$, with $V_{\mathrm{ion}}$ the ionized gas
volume plus the central cavity, and the shell-averaged
$x_{\mathrm{HII}} = 0.5$ crossing that the tables carry as
$R_{\mathrm{out}}$ are $1.464$ and $1.464\times10^{19}$~cm at HII40 against
a published 1.45--$1.47\times10^{19}$, and $8.983$ and
$8.988\times10^{18}$~cm at HII20 against a published
8.7--$9.01\times10^{18}$: the two definitions of the outer boundary are not
identical, but they locate the same front, and both fall inside the
published range at both benchmarks, consistent with a recombination budget
of $0.98$--$1.00\,\QH$.  The temperature panels
reproduce the published shape --- a shallow inward decline followed by a
steep rise at the front, where the surviving radiation field is hardest.

\emph{Grid resolution.}  The two three-dimensional codes of the comparison
state what they used, and it is worth placing \MoCHII{} beside them, because
a Monte Carlo photoionization calculation on a Cartesian grid has a
resolution as well as a photon budget.  \citet{ercolano2003} ran the
published \textsc{Mocassin} benchmark columns on $13\times13\times13$ cells
with the source in a \emph{corner}, so that the grid covers one octant of the
spherically symmetric model; in their words, ``although the number of
sampling points along each orthogonal axis is only 13, this is the
equivalent of a one-dimensional code with 273 radial points.''
\citet{wood2004} used $65^3$ cells over the full box with the star at the
center, which they describe as matching a $33^3$ \textsc{Mocassin} octant,
i.e.\ a cell width of ``3\% of the radius,'' and $10^8$ photon packets.  The
\MoCHII{} runs here use a uniform level-6 octree, $64^3 = 262\,144$ leaves
over the full box with the source at the center, and $2^{25}$ packets
in the ionizing band plus $2^{23}$ between 4.78 and 13.6~eV.
\MoCHII{} and \citet{wood2004} are therefore at essentially the same
resolution --- $5.0\times10^{17}$ (over the box extended to 5.2~pc) against
$4.4\times10^{17}$~cm on a side at HII40 --- and both are more than twice as
fine per axis as the octant grid behind the published \textsc{Mocassin}
column.  We did not repeat the
benchmark on a deliberately coarsened grid, because the quantity at issue is
already resolution-converged in the direction that matters: the front
re-refinement test of Section~\ref{sec:refine} shows $R_{\mathrm{eff}}$
changing by $+0.012\%$ between a native level-7 grid and a coarser
re-refined equivalent, and \citet{ercolano2003} themselves attribute to grid
coarseness only their inner-edge temperature $T_{\mathrm{inner}}$ --- which
their coarse cells make too \emph{low}, since a cell carries the temperature
of its center rather than of the illuminated face --- and not
$\langle T[N_{\mathrm{p}}N_{\mathrm{e}}]\rangle$ or the line ratios compared
here.  Grid resolution is thus not a candidate for the differences discussed
above, with the exception of the thin partly neutral edge that feeds
[O\,\textsc{i}] 6300+6363, which one cell width does not resolve in any of
the three-dimensional codes.

The radial ionization and temperature structure of both benchmarks is
shown in Appendix~\ref{app:lines} (Figures~\ref{fig:wood40} and
\ref{fig:wood20}), in the panel layout of \citet{wood2004}.


\emph{Reference-code context.}  The C25 and GF columns demonstrate the
intrinsic spread among implementations and atomic-data choices.  For example,
at HII40 [S\,\textsc{ii}] 6716+6731 is 0.304 (C25) and 0.137 (GF), while
[O\,\textsc{iii}] 4363 is 0.0046 and 0.00235.  At HII20 the two
[S\,\textsc{iii}] entries are 0.209 and 0.293 in C25, compared with 0.445
and 0.501 in GF.  The published spread and the independently calculated C25
column therefore provide more informative reference ranges than any single
column.

\emph{Where \MoCHII{} sits among the codes.}  \citet{pequignot2001} and
\citet{ercolano2003} note that the eight codes do not scatter randomly:
GF, HN, and DP agree more tightly with one another --- and, in the
planetary-nebula cases, with \textsc{Mocassin} --- than with TK, PH, and RS,
a coherence they attribute to more recently updated atomic data and to a
more careful treatment of the diffuse radiation field.  A second grouping
runs along the transfer treatment: the codes that solve the transfer exactly
(RR, PH, and \textsc{Mocassin}) give systematically lower temperatures at
the illuminated inner edge, because the radiation field there is softer than
an approximate treatment makes it.  \MoCHII{} belongs with the first group on
the data axis --- exact hydrogenic and Verner cross sections, Badnell 2023
recombination, CHIANTI v11 collision strengths, the same vintage
\textsc{Cloudy} uses --- and with the second on the transfer axis, since the
analytic direct-field estimator and the explicit diffuse field used for
these benchmarks are an exact treatment in three dimensions; its inner-edge
temperature at HII40 is correspondingly in the lower half of the published
$T_{\mathrm{inner}}$ range.  Its ionization structure and most of its line
ratios land inside the published spread, with the helium ratio and both
outer radii and $L(\mathrm{H}\beta)$ inside the published ranges; its
weighted temperature matches the \textsc{Cloudy} solution to $0.4\%$ at
both benchmarks.  The cooling treatment that
achieves this, the density suppression of Section~\ref{sec:cooldata}, is
what separates it from the zero-density limit that leaves the two benchmarks
$4\%$ and $9\%$ cooler; the fixed-state comparison quantifies the offset between
the two and shows it to be a property of the cooling function, not of the
atomic data.

%==========================================================================
\section{Solution-driven front refinement}
\label{sec:refine}

A distinguishing capability of the octree grid is that the refinement can
be driven by the ionization solution itself.  At a chosen iteration the
octree is rebuilt: a leaf is flagged as ``front'' when
$\epsilon < \xHI < 1-\epsilon$ or a face neighbor differs by more than a
threshold; the new tree carries the maximum level wherever a cell overlaps
a front-flagged leaf and a base level elsewhere.  The gas state and opacity
map over by position, and every dependent array is recreated.  All ranks
rebuild identically from the shared state, and the old grid's
shared-memory windows are freed immediately after the position sampling, so
repeated re-refinement neither leaks nor holds two grids in memory at
once.

On the Str\"omgren configuration, starting at level 5 and rebuilding to
level 7 at iteration 15 (32\,768 $\to$ 467\,272 leaves), the re-refined
grid gives $R_{\mathrm{eff}} = 3.0562$~pc against a native uniform level-7
grid's 3.0559~pc (2\,097\,152 leaves) --- agreement to $+0.012\%$ with
$4.5\times$ fewer cells, and both improve on level 6 ($+0.68\%$ and
$+0.67\%$ against its $+1.01\%$ from the 1D reference), as resolution
convergence expects (Figure~\ref{fig:g4refine}).  The front
profiles overlie exactly.

\begin{figure}[t]
\centering
\includegraphics[width=\columnwidth]{figures/g4_refine_check.pdf}
% source: python/make_g4_refine.py, from results/continuous_energy/m10_drift_test/g1_g4/
%   (inputs tests/g4_refine/)
\caption{Re-refinement test: the front profiles of the re-refined
$5\to7$ octree and the native uniform level-7 grid overlie exactly, at
$4.5\times$ fewer cells.}
\label{fig:g4refine}
\end{figure}

As a demonstration on a realistic geometry, \MoCHII{} was run on an
Illustris-TNG 60-kpc cutout (density and neutral-fraction columns from the
simulation), with a $\QH = 10^{53}\,\mathrm{s^{-1}}$, 40-kK source at the
density peak and thermal balance, metals, and re-refinement enabled
(Figure~\ref{fig:tng}).  Dense clumps cast sharp neutral shadows, the
circumgalactic gas sits at $\xHI \sim 10^{-4}$ with an $\nelec$-weighted mean
$\Te = 7502$~K, and the refinement (32\,768 $\to$ 617\,135 leaves) tracks
the \emph{distributed} ionization topology of a turbulent medium --- clump
skins and partially ionized gas rather than a single thin shell.  This is
the regime for which the volume-integrated convergence criterion
(Section~\ref{sec:iterate}) was designed.

\begin{figure*}[t]
\centering
\includegraphics[width=0.7\textwidth]{figures/g4_tng_maps.pdf}
% source: python/make_g4_tng.py, from tests/g4_tng/
\caption{\MoCHII{} on an Illustris-TNG cutout: gas density, converged
neutral fraction (note the clump shadows), electron temperature, and the
re-refined cell size, in the source plane.}
\label{fig:tng}
\end{figure*}

%==========================================================================
\section{Dust and PAH coupling}
\label{sec:dust}

Because \MoCHII{} is built on a dust radiative-transfer engine, the dust
physics is a first-class part of the transport.  The coupling has four
parts.

\emph{Absorption competition.}  A packet of energy $E$ sees a dust
absorption opacity $\kappa_{\mathrm{d}} = \rho\kappa\,s_{\mathrm{abs}}(E)$
with $s_{\mathrm{abs}}(E) = (1-a(E))\,C_{\mathrm{ext}}(E)/
C_{\mathrm{ext}}(\lambda_{\mathrm{ref}})$, evaluated at its own energy.  In the default path these
quantities come from the shared astrodust grain model and its EUV Q table
(1762 wavelengths from $1.0\times10^{-4}$ to
$3.981\times10^{4}\,\mu$m), so the grains that absorb the ionizing field
are the same grains used by the SEDust emission solve.  The dusty benchmark
runs below explicitly select a WD01/D03 $R_V = 3.1$ extinction table
\citep{weingartner2001,draine2003} covering the EUV, where
$C_{\mathrm{abs}}/\mathrm{H} = 1.75\times10^{-21}$ at 13.6~eV falls to
$3.96\times10^{-22}\,\mathrm{cm^2}$ at 100~eV.  On the dusty Str\"omgren
test, against a 1D reference carrying identical dust absorption, the full
MW dust gives $R_{\mathrm{eff}} = 2.212$~pc versus the reference 2.197~pc
($+0.70\%$, against $+1.02\%$ for the dust-free run on the same level-6
grid, the finite width of the front on its 0.125~pc cells), reproducing the
classic dusty reduction $R/R_0 = 0.721$ (0.724 in 1D;
Figure~\ref{fig:duststrom}); dust absorbs 61\% of the band luminosity
there.

\emph{Scattering.}  With scattering enabled the scattering opacity joins
the extinction (in-band Henyey--Greenstein asymmetry $g = 0.67$--$0.99$,
strongly forward), and interactions are sampled by a forced-first-flight
scheme after the analytic direct walk, with an albedo weight at each
interaction and Russian roulette.  The full-scattering Str\"omgren radius,
2.203~pc, is within $0.4\%$ of the absorption-only 2.212~pc, and well above
the 2.057~pc of the 1D reference that counts every scattering as an
absorption: the in-band
dust is strongly forward-throwing, so a scattered photon is barely deflected
and scattering acts almost as pure transmission, exactly as expected.

\emph{Live ionization-dependent dust with a PAH split.}  Following
\citet{laursen2009}, the dust content is traced by the metallicity-scaled
hydrogen, with neutral gas carrying the full dust column and ionized gas
only a survival fraction $f_{\mathrm{ion}}$ of it (grains are eroded where
the hydrogen is ionized).  \MoCHII{} evaluates this on the \emph{computed}
neutral fraction and refreshes it every iteration, and it splits the grains
into a large-grain share $1 - f_{\mathrm{PAH}}$ and a PAH share
$f_{\mathrm{PAH}}$ with independent ionized-gas survival, so the dust
opacity of a leaf is
\begin{multline}
  \kappa_{\mathrm{d}} \propto \frac{Z}{Z_{\mathrm{ref}}}\,\nH
  \big[(1 - f_{\mathrm{PAH}})(\xHI + f_{\mathrm{ion}}\,x_{\mathrm{HII}}) \\
  + f_{\mathrm{PAH}}(\xHI + f_{\mathrm{ion,PAH}}\,x_{\mathrm{HII}})\big],
\end{multline}%
with $x_{\mathrm{HII}}$ the stored ionized fraction ($1 - \xHI$) and $f_{\mathrm{ion,PAH}} = 0$ by
default (PAHs destroyed in ionized gas).  Because this dust opacity is
rebuilt from the ionization state each iteration while that state is in
turn set by the dust-attenuated field, the dust--ionization feedback loop
is solved self-consistently; it converges to a near dust-free interior
($R/R_0 = 0.995$).

Figure~\ref{fig:duststrom} contrasts the three treatments on the same
Str\"omgren configuration.  Static full MW dust competes with hydrogen for
the ionizing photons throughout the volume, so the ionization front is
pushed inward to $R_{\mathrm{eff}} = 2.21$~pc, the classic dust-limited
size ($R/R_0 = 0.721$).  Under \code{laursen09\_live}, by contrast, the
interior dust vanishes together with the neutral fraction it is tied to ---
once the gas is ionized it carries essentially no dust to absorb the
ionizing field --- so the front profile overlies the dust-free case almost
exactly ($R_{\mathrm{eff}} = 3.05$~pc, $R/R_0 = 0.995$).  The physical
distinction is that only the neutral gas beyond the front and its thin skin
retain dust, which is where the dust column that matters for the escaping
radiation actually sits.

\begin{figure}[t]
\centering
\includegraphics[width=\columnwidth]{figures/dust_stromgren.pdf}
% source: python/make_dust_stromgren.py, from
%   tests/g1_stromgren/g1_stromgren_rates.h5 (dust-free; these three rerun 2026-09-26),
%   tests/d_dusty/d_dust_full_rates.h5 (full MW dust),
%   tests/d_dusty/d_dust_l09_rates.h5 (laursen09_live).
\caption{H\,\textsc{ii} region size versus dust treatment: the radial
neutral-fraction profile $\xHI(r)$ for the dust-free, full MW dust, and
\code{laursen09\_live} runs of the same Str\"omgren configuration
($\nH = 100\,\mathrm{cm^{-3}}$, $\QH = 10^{49}\,\mathrm{s^{-1}}$,
$T_\star = 4\times10^4$~K, case B).  Dotted vertical lines mark each
effective ionized radius $R_{\mathrm{eff}}$.  Full MW dust pushes the
ionization front inward to $R/R_0 = 0.721$, while the live
ionization-dependent dust vanishes in the ionized interior and recovers the
dust-free size ($R/R_0 = 0.995$).}
\label{fig:duststrom}
\end{figure}

\emph{Grain heating, equilibrium temperature, and IR emission.}  The
transported grain heating drives the equilibrium dust temperature by
balancing the absorbed power against $4\pi\int s_{\mathrm{abs}}(\nu)
B_\nu(T)\,d\nu$; with the band widened to 6~eV, $T_{\mathrm{d}}$ runs from
87~K at 0.3~pc to 28~K at the sphere edge (the canonical \HII-region
gradient), and the IR spectrum integral closes on the absorbed power to
$L_{\mathrm{IR}}/L_{\mathrm{abs}} = 1.000$.  A full stochastic dust
emission solve (the \MoCafe{} SEDust module; astrodust, DL07, and Zubko
grain models with PAH bands, after \citealt{drainel12007}) is available at
output time: it uses the size-resolved stochastic-heating spectrum as a
\emph{shape} normalized to the transported absorbed power, so the
normalization is exact by construction.  The reradiated infrared is then
itself transported: packets are launched from each cell in proportion to
its emitted power, with the wavelength drawn continuously from that
cell's emission spectrum, and carried through the same grain opacity that
produced the emission, so what the grains reabsorb returns to their
heating.  The transport confirms rather than assumes that an \HII{}
region is thin to its own infrared --- the dusty Str\"omgren model
reabsorbs 0.6--0.9\% of it and converges in two passes --- and it
supplies the escaping infrared spectrum directly.  With the live PAH share
(Figure~\ref{fig:pah}), the PAH features between 3.3 and 12.7\,$\mu$m drop to 0.13--0.24 of the
PAH-everywhere reference and the projected 8\,$\mu$m map turns from
centrally peaked into a \emph{limb-brightened annulus} just inside the
ionization front --- the classic Spitzer/JWST bubble morphology.

\begin{figure*}[t]
\centering
\includegraphics[width=0.95\textwidth]{figures/pahlive_maps.pdf}
% source: python/make_dust_maps.py, from
%   tests/d_dusty/dustemis_smoke_dustemis.h5 (PAHs everywhere) and
%   tests/d_dusty/pahlive_smoke_dustemis.h5 (live PAH survival).
%   Both smokes rerun with the binary of 2026-09-27 and the figure
%   regenerated the same day.
\caption{Stochastic dust emission at 8\,$\mu$m with the
ionization-dependent PAH share.  The first two panels are two-dimensional
projection maps sharing one brightness scale: the band emissivity, which
\MoCHII{} writes as the power radiated per unit volume and per unit
wavelength into $4\pi$~sr, integrated along the line of sight and divided by
$4\pi$, so the panels show the specific intensity $I_\lambda$ of an
infrared-thin nebula.  Left: PAHs surviving
everywhere, giving a centrally peaked image.  Middle: PAHs destroyed in the
ionized interior ($x_{\mathrm{HII}}$-weighted survival).  Right: the
azimuthally averaged radial profiles of both maps.  Destroying the PAHs
removes the central peak and leaves a broad limb-brightened annulus just
inside the ionization front at $r \simeq 2.3$~pc, visible in the profile as
a local maximum a factor 1.6 above the depressed interior.}
\label{fig:pah}
\end{figure*}

\begin{figure*}[t]
\centering
\includegraphics[width=0.95\textwidth]{figures/dustemis_maps.pdf}
% source: python/make_dust_maps.py, from
%   tests/d_dusty/dustemis_smoke_dustemis.h5.
%   Rerun with the binary of 2026-09-27; figure regenerated the same day.
\caption{Dust emission at 8 and 160\,$\mu$m from the stochastic SEDust
solve, as two-dimensional projection maps of the specific intensity
$I_\lambda$ (left, middle) with their azimuthally averaged radial profiles
(right).  The band emissivity is written as the power radiated per unit
volume and per unit wavelength into $4\pi$~sr, so $I_\lambda$ is that
emissivity integrated along the line of sight and divided by $4\pi$.  The
short
wavelength is centrally peaked (warm inner dust), falling by 3.5~dex across
the sphere, while the far-infrared profile is far flatter (1.45~dex), the
projected signature of the dust temperature gradient.}
\label{fig:dustemis}
\end{figure*}

%==========================================================================
\section{The photodissociation-region skin}
\label{sec:pdr}

With the FUV band extension and three switches, \MoCHII{} models the
atomic photodissociation-region (PDR) skin just outside the ionization
front.  The switches feed the metal-cascade electrons into the $\nelec$
closure and add metal photoheating, with its recombination and
collisional-ionization losses, to the thermal balance (both on by default),
and add grain photoelectric heating and grain recombination cooling with
the \citet{bakes1994} fits driven by the local FUV field $G_0$ (off by
default).  The coolant is the ground-term fine structure of
C\,\textsc{i}, C\,\textsc{ii} and O\,\textsc{i} excited by atomic hydrogen
as well as by electrons, solved in one statistical equilibrium
(Section~\ref{sec:cooldata}; \citealt{abrahamsson2007,barinovs2005}).  It
is not a switch but part of the metal cooling whenever metals are carried,
and it is essential here: with $n_{\mathrm{HI}}/\nelec \sim 10^3$--$10^4$ in
the PDR the hydrogen collisions set the level populations, and the
[O\,\textsc{i}] 63 and 145\,$\mu$m, [C\,\textsc{i}] 609 and 370\,$\mu$m and
[C\,\textsc{ii}] 158\,$\mu$m emissivities written for the zone come from the
same populations that cool it.  Two guards keep the
photoelectric heating physical (evaluated at $\min(T, 10^4\,\mathrm{K})$ so
the fits stay in range, and weighted by $\xHI$ so ionized gas, whose grains
are EUV-charged and whose energy is already routed to the IR, is not
overheated).  Both are necessary: without them the fits are evaluated
outside their validity range and the zone runs away to a false
12--35~kK.

Figure~\ref{fig:pdr} shows the converged radial structure of a PDR test
(a $\nH = 100$ octree sphere with a central source and all three switches
on).  The interior is fully ionized with a sharp front; beyond the front,
the C\,\textsc{ii} fraction is unity, the electron density falls to the
carbon-dominated value $\nelec \simeq A_{\mathrm{C}}\nH$, and the electron
temperature settles to a genuine PDR equilibrium (a gradient from
$\sim$250~K just beyond the front to $\sim$170~K at the outer edge, not
pinned at the temperature floor).  Against the code before the
atomic-data update of Section~\ref{sec:atomic}, whose two-level hydrogen
impact entered the cooling but not the written lines, the zone emitted $1.98$
times the [C\,\textsc{ii}] 158\,$\mu$m and $39.5$ times the
[O\,\textsc{i}] 63\,$\mu$m luminosity, and the neutral shell was cooler (it
ran from $\sim$400 to $\sim$200~K before); those factors were measured
before the \HeI{} and \HeII{} on-the-spot photons
(Section~\ref{sec:diffuse}) and the other current defaults were in place.
The run of Figure~\ref{fig:pdr} carries them, together with the hydrogenic
$\alpha_B$ of Section~\ref{sec:recomb} and the helium--hydrogen charge
exchange of Section~\ref{sec:ci}, and against the same model converged
without any of these it has $5.8\%$ more [C\,\textsc{ii}] 158\,$\mu$m
and $3.5\%$ less [O\,\textsc{i}] 63\,$\mu$m luminosity, $6.9\%$ more
H$\beta$, and the same neutral-zone temperature profile.  This is the
atomic PDR skin only: there is no H$_2$/CO chemistry or line transfer, so
the PDR temperatures are first-order estimates, adequate for the ion
structure and rough [C\,\textsc{ii}] emissivities.

\begin{figure}[t]
\centering
\includegraphics[width=\columnwidth]{figures/pdr_profiles.pdf}
% source: python/make_pdr_profiles.py, from
%   results/continuous_energy/m11_hecx_solver/pdr_L5_converged_rates.h5 (tests/pdr/pdr_L5.in, Sobol launch,
%   restarted from m5_converged with the 2026-09-25 defaults, the hydrogenic alpha_B and
%   the stopping test of 2026-09-26, the He+ + H charge exchange and
%   hydrogen-balance solve of 2026-09-27)
\caption{Radial structure of the PDR test: the neutral fraction, electron
temperature, and electron density as a function of radius.  Inside the
front the gas is fully ionized; beyond it the electron density drops to the
carbon-dominated value and the temperature settles to a PDR equilibrium
above the floor.}
\label{fig:pdr}
\end{figure}

%==========================================================================
\section{Summary and outlook}
\label{sec:summary}

\MoCHII{} is a three-dimensional Monte Carlo photoionization code for dusty
nebulae on adaptive octree grids, built by adding a self-consistent
gas-physics layer to a validated dust radiative-transfer engine.  Its
design choices --- octree AMR with solution-driven front refinement, a
zero-variance analytic direct-field estimator with a dedicated Cartesian
DDA backend, an explicit Monte Carlo diffuse field with an on-the-spot
option, a
data-driven trace-metal registry with modern atomic data, first-class dust
absorption/scattering and equilibrium/stochastic emission with PAH physics,
and peel-off synthetic imaging --- distinguish it from the uniform-grid and
1D codes.  It reproduces analytic attenuation to $+0.08\%$, the Str\"omgren
sphere to $+0.69\%$ in radius, and a thermal H/He nebula to $0.12\%$.  On the
Lexington 2000 \HII{} region benchmarks, with continuously sampled Planck
photon energies and cross sections evaluated at each sampled energy, and
with line cooling taken from the level populations at the local electron
density, the
$n_{\mathrm{p}}n_{\mathrm{e}}$-weighted temperature is 8180~K at HII40 and
6940~K at HII20, matching the \textsc{Cloudy} c25.00 solution of the same
two models (8210 and 6910~K) to $0.4\%$, and the HII40 \HeI{} front lies
within $0.8\%$ in radius of \textsc{Cloudy}'s; both models are
radiation-bounded, so both outer radii are predicted and fall inside the
published ranges, as does $L(\mathrm{H}\beta)$.  These are the converged
solutions: a stopping test on the last iteration's change alone had ended
the runs while the fronts were still advancing, and the test now bounds the
drift still to come, for helium as well as hydrogen, and accepts the Monte
Carlo floor only when the ionized volume shows no net drift.  The
line-cooling data are checked independently of the temperature by evaluating
the \MoCHII{} cooling terms on a fixed state taken from a \textsc{Cloudy}
c25.00 cooling decomposition: channel by channel the two compilations agree,
and the summed line cooling agrees to $1\%$.  The zero-density limit of the
same cooling data, retained as an option, overestimates the line cooling
by $12\%$ at $\nelec = 110\,\mathrm{cm^{-3}}$, all of it in the far-infrared
fine-structure lines, and places both
benchmarks $4$--$9\%$ cooler (7877 and 6282~K, inside the published range at
HII40 and below it at HII20); the density suppression that removes it is the
default.
The \textsc{Cloudy} c25.00 column provides a uniform reference for both
tables and has the same S\,\textsc{ii}/S\,\textsc{iii} ordering as
\MoCHII{}.
The solution-driven
re-refinement reproduces a native high-resolution grid to $+0.012\%$ with
$4.5\times$ fewer cells, the dust coupling closes energy exactly and
reproduces the limb-brightened Spitzer/JWST PAH morphology, and the
atomic PDR skin
recovers the carbon-dominated electron density and a physical PDR
temperature.

The current scope is restricted to stellar-ionizing spectra below 100~eV and
an atomic PDR treatment.  Molecular chemistry and higher ionization stages
are not included; inner-shell photoabsorption is included, but the Auger
ejection that follows it for the Mg, Ca and Fe vacancies of
Section~\ref{sec:xsec} is not.  Extensions to plane-parallel illumination, harder
spectra, and resonant-scattering transfer of the \HeI{} 10833~\AA\ line
through the coupling to \textsc{LaRT} will broaden the range of
applications to planetary nebulae, active galactic nuclei, and illuminated
interstellar and circumgalactic gas.

\begin{acknowledgments}
K.-I.\ Seon was supported by the Korea Astronomy and Space Science
Institute under the R\&D program (Project No.\ 2026-1-831-02) supervised by
the Korea AeroSpace Administration.
\end{acknowledgments}

\software{MoCHII (\url{https://github.com/seoncafe/MoCHII}),
          MoCafe (\url{https://github.com/seoncafe/MoCafe}),
          CHIANTI \citep{dere1997,delzanna_chianti},
          MOCASSIN \citep{ercolano2003},
          Cloudy \citep{cloudy2023},
          astropy,
          matplotlib,
          NumPy}

\appendix

\section{Fitted atomic-data coefficients}
\label{app:atomic}

The atomic data are fitted offline and evaluated online
(Section~\ref{sec:cooldata}).  Table~\ref{tab:tier1err} gives the quality of
the line-cooling fits of every ion and Table~\ref{tab:tier2} summarizes the
$n$-level diagnostic models; it is generated mechanically from the data
files by \code{tables/make\_atomic\_tables.py}.  The coefficients
themselves --- the amplitudes and excitation temperatures of the
line-cooling fits, and the Burgess--Tully Chebyshev $\Upsilon(T)$ sets of
the $n$-level models, several for each transition --- are distributed with
the code rather than reproduced here.\footnote{The fitted line-cooling
coefficients of every ion are provided as text in
\code{data/atomic/cooling\_fit\_parameters.txt} (and, one file for each ion,
in \code{data/atomic/cooling\_<ion>.txt}), with the fit form, units, fit
range, and maximum fit error documented in \code{data/atomic/README.md};
the $n$-level models are in \code{data/atomic/nlevel\_<ion>.txt}.  Every
file carries a provenance header.}

\clearpage
\input{tables/tier1_maxerr}

\input{tables/tier2_nlevel}

\section{Full Lexington benchmark line lists}
\label{app:lines}

\begin{figure*}[t]
\centering
\includegraphics[width=\textwidth]{figures/wood_hii40.pdf}
% source: results/continuous_energy/m11_hecx_solver/hii40_converged_rates.h5
% and the Cloudy c25.00 save files
% tests/cloudy_c25/hii40_c25.{ovr,ele_h,ele_he,ele_c,ele_n,ele_o,ele_ne},
% via python/make_wood_benchmark.py
\caption{Lexington HII40 benchmark: radial ionization and temperature
structure of the \MoCHII{} solution, in the panel layout of
\citet{wood2004}, their Figure 1.  Each point is one octree leaf inside the
prescribed shell, thinned at random for legibility; the temperature panel
shows the leaves that are less than 95\% neutral, the convention of
\citet{wood2004}.  Nitrogen stops at N$^{+2}$ because the registry carries
three nitrogen stages for this model.  This is the same run as the
\MoCHII{} column of Table~\ref{tab:wood40}, with continuous Planck-sampled
stellar photon energies and cross sections evaluated at the sampled energy;
no binned radiation field enters it.  Black lines are
\textsc{Cloudy}~c25.00 on the same
model --- the run behind the C25 column of Table~\ref{tab:wood40}, whose
input is the unmodified c25.00 test-suite deck
\code{tsuite/auto/hii\_paris.in} --- drawn as a line as
\citet{wood2004} draw \textsc{Cloudy} in their own figures.  The
\textsc{Cloudy} profile is cut at the outer radius of the \MoCHII{} shell so
that both codes cover one volume, and the 95\% cut is applied to its zones
as well.  In the temperature panel the \MoCHII{} points run below the
\textsc{Cloudy} curve through the interior and cross it at the front.}
\label{fig:wood40}
\end{figure*}

\begin{figure*}[t]
\centering
\includegraphics[width=\textwidth]{figures/wood_hii20.pdf}
% source: results/continuous_energy/m11_hecx_solver/hii20_converged_rates.h5
% and the Cloudy c25.00 save files
% tests/cloudy_c25/hii20_c25.{ovr,ele_h,ele_he,ele_c,ele_n,ele_o,ele_ne},
% via python/make_wood_benchmark.py
\caption{Lexington HII20 benchmark, in the panel layout of
\citet{wood2004}, their Figure 2; conventions as in
Figure~\ref{fig:wood40}.  Ne$^{+2}$ is confined within 1.4~pc, and
C$^{+3}$ stays below the plotted range and carries no series.  This is the
same run as the \MoCHII{} column of Table~\ref{tab:wood20}.  At 20~kK the
helium-ionizing photons are so few that helium is neutral beyond 1.4~pc.
The continuous-energy treatment nonetheless preserves the weak C$^{+2}$
tail outside that helium transition.  Black lines are \textsc{Cloudy}~c25.00
on the same model, from the
unmodified c25.00 deck \code{tsuite/auto/hii\_coolstar.in} behind the C25
column of Table~\ref{tab:wood20}.  \textsc{Cloudy} stops on its own
criterion at 2.91~pc, short of the outer radius of the \MoCHII{} shell, so
its curves end there, and its final near-vertical fall is the recombination
front itself, resolved in zones far finer than a \MoCHII{} cell.}
\label{fig:wood20}
\end{figure*}

Tables~\ref{tab:wood40} and~\ref{tab:wood20} give the complete Lexington
HII40 and HII20 benchmark line lists discussed in Section~\ref{sec:hii}, in
the form used by \citet{wood2004}: the eight published codes of the
\citet{ercolano2003} compilation, the \citet{wood2004} Monte Carlo code
(WME), a \textsc{Cloudy}~c25.00 calculation (C25), and the \MoCHII{} column.
Line intensities are relative to H$\beta$.

% Tables tab:wood40 / tab:wood20: generated by tables/make_wood_tables.py.
% Published columns transcribed from the LaTeX source of Wood et al. (2004)
% (arXiv:astro-ph/0311584, kwood.tex) and checked cell by cell against it;
% the MoCHII column is computed from the continuous-energy benchmark lines
% in results/continuous_energy/m11_hecx_solver/.
\clearpage
\input{tables/wood_hii40}
\input{tables/wood_hii20}

\bibliographystyle{aasjournalv7.1}
\bibliography{mochii}

% Backup copy of the pre-restructuring manuscript.
\end{document}
